% Minimal preamble for Open Logic Project sources (pinned commit 1e960beff9ed7835bf3e3f1335e21af3439cd107).
%
% OLP's own master preamble cannot run on the project host: TeX Live 2025 there lacks
% prooftree.sty, ntheorem.sty, extarrows.sty, mathdots.sty, isomath.sty and subfiles.sty.
% Following the project's own rule ("if the author preamble needs packages the toolchain lacks,
% do not invent macros -- use a minimal working preamble"), this file declares, verbatim in
% spirit and numbering behaviour, exactly the theorem-like environments OLP defines in
% sty/open-logic-book-envs.sty (thm/lem/prop/cor/defn/ex/prob/rem/note/case/axiom/conv),
% plus the shorthand accents OLP uses (`!!a`, `!!b`, `!!{...}`).
%
% Author text is untouched; only the packaging is ours.
\documentclass[10pt,a4paper]{article}

\usepackage{amsmath}
\usepackage{amssymb}
\usepackage{amsthm}
\usepackage{graphicx}
\usepackage{hyperref}
\usepackage{enumitem}
\usepackage[framemethod=default]{mdframed}

% ---- OLP shorthand accents (from sty/open-logic.sty) ----
% NB: \! is taken by LaTeX (negative thin space), so we only define the !!a / !!b families
% and the \olup helper used by them.
\newcommand*{\olup}[1]{\MakeUppercase{#1}}
\def\!!!a#1{\olup{#1}}
\def\!a#1{\olup{#1}}
\def\!b#1{\olup{#1}}

% ---- theorem-like environments, same family numbering as OLP ----
\theoremstyle{plain}
\newtheorem{thm}{Theorem}[section]
\newtheorem{lem}[thm]{Lemma}
\newtheorem{prop}[thm]{Proposition}
\newtheorem{cor}[thm]{Corollary}
\newtheorem{axiom}{Axiom}

\theoremstyle{definition}
\newtheorem{defn}[thm]{Definition}
\newtheorem{ex}[thm]{Example}
\newtheorem{prob}{Problem}[section]
\newtheorem{probd}{Problem}
\newtheorem{conv}[thm]{Convention}

\theoremstyle{remark}
\newtheorem{rem}[thm]{Remark}
\newtheorem{note}[thm]{Note}
\newtheorem{case}[thm]{Case}

% ---- environments OLP defines as wrappers ----
\newenvironment{probtag}[1]{\begin{prob}}{\end{prob}}
\newenvironment{explain}{}{}
\newenvironment{explain*}{}{}
\newenvironment{editorial}{}{}
\newenvironment{digress}{}{}
\newenvironment{digress*}{}{}
\newenvironment{history}{}{}
\newenvironment{reference}{}{}
\newcommand*{\olfileid}[3]{}
\newcommand*{\olsection}[1]{\section*{#1}}
\newcommand*{\olchapter}[3]{}
\newcommand*{\olpart}[2]{}
\newcommand*{\olimport}[1]{}
\newcommand*{\olinclude}[1]{}
\newcommand*{\olpath}[1]{}
\newcommand*{\ollangid}[1]{}
\newcommand*{\problemtag}[1]{}
\newcommand*{\problemauthor}[1]{}
\newcommand*{\ollabel}[1]{\label{#1}}
% \olref takes 0-3 OPTIONAL arguments plus one mandatory label:
%     \olref{lab}  |  \olref[from]{lab}  |  \olref[from][to]{lab}  |  \olref[a][b][c]{lab}
% A single-argument definition silently swallows "[sth]" and emits a broken reference
% (measured on the host: the text came out as "??iso]x"). The chain below consumes the
% brackets and cross-references only the mandatory argument.
\makeatletter
\newcommand{\olref}{\@ifnextchar[{\olref@i}{\olref@final}}
\def\olref@i[#1]{\@ifnextchar[{\olref@ii}{\olref@final}}
\def\olref@ii[#1]{\@ifnextchar[{\olref@iii}{\olref@final}}
\def\olref@iii[#1]{\olref@final}
\def\olref@final#1{\ref{#1}}
\makeatother
\newcommand*{\olnewpage}{\clearpage}

% ===================== OLP compatibility layer (1) =====================

\makeatletter

% ---------------- proof trees (OLP's prooftree is a bussproofs derivative) ----------------
\providecommand{\AxiomC}[1]{}
\providecommand{\UnaryInfC}[1]{}
\providecommand{\BinaryInfC}[1]{}
\providecommand{\TrinaryInfC}[1]{}
\providecommand{\QuaternaryInfC}[1]{}
\providecommand{\QuinaryInfC}[1]{}
\providecommand{\RightLabel}[1]{}
\providecommand{\LeftLabel}[1]{}
\providecommand{\doubleLine}{}
\providecommand{\noLine}{}
\providecommand{\DisplayProof}{}
\providecommand{\fCenter}{}
\newenvironment{prooftree}{}{}
\newenvironment{ltprooftree}{}{}
\newenvironment{oltableau}{}{}

% ---------------- logical notation ----------------
\providecommand{\lif}{\mathbin{\rightarrow}}
\providecommand{\lforall}{\mathord{\forall}}
\providecommand{\lexists}{\mathord{\exists}}
\providecommand{\ltrue}{\top}
\providecommand{\lfalse}{\bot}
\providecommand{\Ltrue}{\top}
\providecommand{\Lfalse}{\bot}
\providecommand{\Sequent}{\mathrel{\Rightarrow}}
\providecommand{\Proves}{\mathrel{\vdash}}
\providecommand{\Entails}{\mathrel{\models}}
\providecommand{\Subst}[3]{#1[#2/#3]}
\providecommand{\sFmla}[2]{#1(#2)}
\providecommand{\pFmla}[2]{#1(#2)}
\providecommand{\formula}{\relax}
\providecommand{\True}{\top}
\providecommand{\False}{\bot}
\providecommand{\Args}[1]{#1}
\providecommand{\Ts}{\ensuremath{\top}}
\providecommand{\Fi}{\ensuremath{\bot}}
\providecommand{\lfalsum}{\bot}

% ---------------- objects, structures, metatheory ----------------
\providecommand{\Obj}{}
\providecommand{\Struct}{\ensuremath{\mathfrak{A}}}
\providecommand{\Domain}{\ensuremath{|\mathfrak{A}|}}
\providecommand{\Assign}{\ensuremath{s}}
\providecommand{\Sat}[2]{\ensuremath{#1 \models #2}}
\providecommand{\mSat}[3]{\ensuremath{#1 \models #2[#3]}}
\providecommand{\mModel}[2]{\ensuremath{#1 \models #2}}
\providecommand{\Lang}{\ensuremath{\mathcal{L}}}
\providecommand{\Th}{\ensuremath{\mathrm{Th}}}
\providecommand{\Log}{\ensuremath{\mathrm{L}}}
\providecommand{\tuple}[1]{\langle #1 \rangle}
\providecommand{\num}[1]{\underline{#1}}
\providecommand{\cfind}[1]{\varphi_{#1}}
\providecommand{\Nat}{\ensuremath{\mathbb{N}}}
\providecommand{\Mult}{\ensuremath{\mathrm{Mult}}}
\providecommand{\Add}{\ensuremath{\mathrm{Add}}}
\providecommand{\Exp}{\ensuremath{\mathrm{Exp}}}
\providecommand{\Zero}{\ensuremath{0}}
\providecommand{\Succ}{\ensuremath{\mathrm{S}}}
\providecommand{\lambd}{\ensuremath{\lambda}}
\providecommand{\bredpar}{\mathrel{\rightarrow_\beta}}
\providecommand{\bcd}[1]{\ensuremath{\ulcorner #1 \urcorner}}
\providecommand{\Set}[1]{\ensuremath{\{#1\}}}
\providecommand{\indcase}[1]{\par\noindent\emph{Case #1}.\quad}
\providecommand{\indfrm}[1]{\par\noindent\emph{Form #1}.\quad}
\providecommand{\iftag}[1]{#1}
\providecommand{\tagitem}[1]{\item[#1]}
\providecommand{\Elim}[1]{}
\providecommand{\Intro}[1]{}
\providecommand{\LeftR}[1]{}
\providecommand{\RightR}[1]{}
\providecommand{\TAss}[1]{}
\providecommand{\TRule}[1]{}

\makeatother
\providecommand{\olpath}{.}

% ===================== OLP compatibility layer (2) =====================
% --- second compatibility block: notation and shorthands used by the frozen OLP units ---
\providecommand{\CutCS}{\mathrm{Cut}}
\providecommand{\LeftSubproofLabel}{\relax}
\providecommand{\RightSubproofLabel}{\relax}
\providecommand{\maeh}{\relax}
\providecommand{\cheight}{\relax}
\providecommand{\ande}{\mathrm{and}_E}
\providecommand{\andi}{\mathrm{and}_I}
\providecommand{\ore}{\mathrm{or}_E}
\providecommand{\ori}{\mathrm{or}_I}
\providecommand{\typeof}{\mathord{:}}
\providecommand{\pValue}[1]{#1}
\providecommand{\Sent}{\mathrm{Sent}}
\providecommand{\cardplus}{\mathbin{+}}
\providecommand{\cardtimes}{\mathbin{\cdot}}
\providecommand{\cardexpo}{\mathbin{\hat{\ }}}
\providecommand{\cardnless}{\not<}
\providecommand{\cardle}{\leq}
\providecommand{\funfromto}[3]{#1\colon #2 \to #3}
\providecommand{\funimage}[1]{\operatorname{im}(#1)}
\providecommand{\fdefined}{\mathord{\downarrow}}
\providecommand{\canonord}{\operatorname{ord}}
\providecommand{\ordtimes}{\mathbin{\cdot}}
\providecommand{\ordeq}[2]{#1 = #2}
\providecommand{\ordneq}[2]{#1 \neq #2}
\providecommand{\ordtype}[1]{\operatorname{ord}(#1)}
\providecommand{\ordsucc}{\mathord{+1}}
\providecommand{\ordunion}{\cup}
\providecommand{\onesphere}{S^1}
\providecommand{\rotationsgroup}{\mathrm{SO}(2)}
\providecommand{\equivrep}[1]{[#1]}
\providecommand{\rlexless}{\prec}
\providecommand{\ZFC}{\mathrm{ZFC}}
\providecommand{\ZF}{\mathrm{ZF}}
\providecommand{\Rat}{\mathbb{Q}}
\providecommand{\Real}{\mathbb{R}}
\providecommand{\PosInt}{\mathbb{Z}^{+}}
\providecommand{\Id}{\mathrm{id}}
\providecommand{\Bin}{\mathrm{Bin}}
\providecommand{\pto}{\rightharpoonup}
\providecommand{\dom}{\operatorname{dom}}
\providecommand{\ran}{\operatorname{ran}}
\providecommand{\len}{\operatorname{len}}
\providecommand{\concat}{\mathbin{{}^\frown}}
\providecommand{\card}[1]{\lvert #1 \rvert}
\providecommand{\Pow}{\mathcal{P}}
\providecommand{\Setabs}[2]{\{\, #1 \mid #2 \,\}}
\providecommand{\comp}{\circ}
\providecommand{\Char}{\chi}
\providecommand{\Prf}{\operatorname{Prf}}
\providecommand{\Prov}{\operatorname{Prov}}
\providecommand{\Proj}{\operatorname{proj}}
\providecommand{\Gn}{\operatorname{gn}}
\providecommand{\gn}{\operatorname{gn}}
\providecommand{\Value}[1]{\operatorname{val}(#1)}
\providecommand{\Deduce}{\mathrm{D}}
\providecommand{\Dual}{\mathrm{d}}
\providecommand{\Ax}{\mathrm{Ax}}
\providecommand{\pSat}[3]{#1 \models #2\,[#3]}
\providecommand{\pAssign}{\sigma}
\providecommand{\bforall}{\forall}
\providecommand{\bexists}{\exists}
\providecommand{\liff}{\mathbin{\leftrightarrow}}
\providecommand{\FV}{\operatorname{FV}}
\providecommand{\UnaryInf}{\ensuremath{\mid}}
\providecommand{\BinaryInf}{\ensuremath{\mid}}
\providecommand{\depth}{\operatorname{depth}}
\providecommand{\oliflabeldef}[3]{#1}
\providecommand{\TMblank}{\sqcup}
\providecommand{\TMstroke}{\mid}
\providecommand{\TMstay}{\mathrm{S}}
\providecommand{\TMright}{\mathrm{R}}
\providecommand{\TMleft}{\mathrm{L}}
\providecommand{\TMendtape}{\triangleright}
\providecommand{\umin}{\mu}
\providecommand{\Sum}{\sum}
\providecommand{\citep}[1]{\cite{#1}}
\providecommand{\citet}[1]{\cite{#1}}
\providecommand{\citeauthor}[1]{\cite{#1}}
\providecommand{\citeyear}[1]{\cite{#1}}

% ===================== OLP compatibility layer (3) =====================
% --- third compatibility block: remaining notation seen in the frozen OLP units ---
\providecommand{\iso}{\cong}
\providecommand{\nszero}{\mathrm{0}}
\providecommand{\nssucc}{\mathrm{S}}
\providecommand{\nsplus}{\mathbin{+}}
\providecommand{\nssless}{<}
\providecommand{\nsless}{<}
\providecommand{\PL}{\mathrm{PL}}
\providecommand{\Nec}{\mathord{\Box}}
\providecommand{\RK}{\mathrm{RK}}
\providecommand{\ST}{\mathrm{ST}}
\providecommand{\SSubst}[3]{#1[#2/#3]}
\providecommand{\subst}[3]{#1[#2/#3]}
\providecommand{\tagfalse}{\mathrm{F}}
\providecommand{\cutr}{\mathrm{Cut}}
\providecommand{\cardneq}{\neq}
\providecommand{\issucc}[1]{\operatorname{succ}(#1)}
\providecommand{\tup}[1]{\langle #1 \rangle}
\providecommand{\Var}{\mathrm{Var}}
\providecommand{\Bew}{\operatorname{Bew}}
\providecommand{\Con}{\operatorname{Con}}
\providecommand{\Cons}{\operatorname{Cons}}
\providecommand{\pDeduce}{\mathrm{D}}
\providecommand{\psat}{\models}
\providecommand{\den}[1]{\llbracket #1 \rrbracket}
\providecommand{\Sem}[2]{\llbracket #1 \rrbracket_{#2}}
\providecommand{\DomainOf}[1]{\lvert #1 \rvert}
\providecommand{\Form}{\mathrm{Form}}
\providecommand{\Terms}{\mathrm{Term}}
\providecommand{\Atom}{\mathrm{Atom}}
\providecommand{\FVars}[1]{\operatorname{FV}(#1)}
\providecommand{\fv}[1]{\operatorname{fv}(#1)}
\providecommand{\tuplep}[1]{\langle #1 \rangle}
\providecommand{\restrict}[2]{#1 \upharpoonright #2}
\providecommand{\sto}{\rightarrow}
\providecommand{\FmlaSet}[1]{\mathrm{Fmla}_{#1}}
\providecommand{\Thm}{\mathrm{Thm}}
\providecommand{\Eval}{\operatorname{val}}
\providecommand{\blank}{\sqcup}
\providecommand{\blankTape}{\sqcup}
\providecommand{\Consis}{\operatorname{Con}}
\providecommand{\drs}{\mathrm{d}}
\providecommand{\VRule}{\mid}
\providecommand{\Bet}{\mathrm{B}}
\providecommand{\Seg}[1]{\mathrm{Seg}_{#1}}
\providecommand{\Enums}{\mathrm{En}}
\providecommand{\Z}{\mathbb{Z}}
\providecommand{\N}{\mathbb{N}}
\providecommand{\abs}[1]{\lvert #1 \rvert}
\providecommand{\floor}[1]{\lfloor #1 \rfloor}
\providecommand{\ceil}[1]{\lceil #1 \rceil}

% ===================== OLP compatibility layer (4) =====================
% --- fourth compatibility block: harvested from the per-candidate compile logs ---
\providecommand{\beredpar}{\mathrel{\rightarrow_\beta}}
\makeatletter
\def\bered{\@ifnextchar[{\bered@i}{\mathrel{\rightarrow_\beta}}}
\def\bered@i[#1]{\mathrel{\rightarrow_\beta}}
\def\bred{\@ifnextchar[{\bred@i}{\mathrel{\rightarrow_\beta}}}
\def\bred@i[#1]{\mathrel{\rightarrow_\beta}}
\makeatother
\providecommand{\beredone}{\mathrel{=_\beta}}
\providecommand{\bredone}{\mathrel{=_\beta}}
\providecommand{\xredone}{\mathrel{=_\beta}}
\providecommand{\aconvone}{\mathrel{=_\alpha}}
\providecommand{\cardless}{<}
\providecommand{\tagrefs}[1]{#1}
\providecommand{\ordplus}{\mathbin{+}}
\providecommand{\funrestrictionto}[2]{#1 \upharpoonright #2}
\providecommand{\fundefined}{\mathord{\uparrow}}
\providecommand{\Trm}{\mathrm{Trm}}
\providecommand{\OProv}{\operatorname{Prov}}
\providecommand{\OPrf}{\operatorname{Prf}}
\providecommand{\aeq}{\mathrel{=}}
\providecommand{\mClass}[1]{[#1]}
\providecommand{\equivclass}[1]{[#1]}
\providecommand{\pheight}{\operatorname{ht}}
\providecommand{\phy}{\operatorname{ht}}
\providecommand{\Cut}{\mathrm{Cut}}
\providecommand{\QuantRank}{\operatorname{qr}}
\providecommand{\elemequiv}{\mathrel{\equiv}}
\providecommand{\PIso}{\mathrm{PIso}}
\providecommand{\Theory}{\mathrm{Th}}
\providecommand{\supstrict}{\subsetneq}
\providecommand{\trcl}{\operatorname{trcl}}
\providecommand{\cref}[1]{\ref{#1}}
\providecommand{\Discharge}{\mathrm{Dis}}
\providecommand{\Weakening}{\mathrm{Wk}}
\providecommand{\ZFminus}{\mathrm{ZF}^{-}}
\makeatletter
\renewcommand{\Dual}{\@ifnextchar[{\Dual@i}{\ifmmode\mathrm{d}\else$\mathrm{d}$\fi}}
\def\Dual@i[#1]{\ifmmode\mathrm{d}[#1]\else$\mathrm{d}[#1]$\fi}
\makeatother
\providecommand{\RightR}{\mathrm{R}}
\providecommand{\LeftR}{\mathrm{L}}
\providecommand{\Intro}{\mathrm{I}}
\providecommand{\Axiom}{\mathrm{Ax}}
\providecommand{\formula}[1]{#1}
\providecommand{\stared}{\ast}
\providecommand{\oset}[2]{\{\, #1 \mid #2 \,\}}
\providecommand{\setof}[2]{\{\, #1 \mid #2 \,\}}
\providecommand{\pow}[1]{\mathcal{P}(#1)}
\providecommand{\ordpair}[2]{\langle #1, #2 \rangle}
\providecommand{\ordtriple}[3]{\langle #1, #2, #3 \rangle}
\providecommand{\union}{\cup}
\providecommand{\inter}{\cap}
\providecommand{\seq}[1]{\langle #1 \rangle}
\providecommand{\SuchThat}{\mid}
\providecommand{\where}{\mid}
\providecommand{\TVar}{\mathrm{TVar}}
\providecommand{\Bool}{\mathrm{Bool}}
\providecommand{\Val}{\operatorname{Val}}
\providecommand{\TL}[1]{\mathrm{TL}}
\providecommand{\xstared}{\ast}

% ===================== OLP compatibility layer (5) =====================
% --- fifth compatibility block: rescues the quarantined OLP candidates ---
\providecommand{\ident}{\operatorname{id}}
\providecommand{\setrank}{\operatorname{rank}}
\providecommand{\cardfont}{\mathfrak{c}}
\providecommand{\Complement}{\overline}
\providecommand{\QR}{\mathrm{QR}}
\providecommand{\becd}{\ensuremath{\ulcorner}}
\providecommand{\disjointsum}{\mathbin{\oplus}}
\providecommand{\cardeq}[2]{#1 = #2}
\providecommand{\xred}{\mathrel{\rightarrow}}
% --- sixth compatibility block: names reported by real compile logs ---
\providecommand{\ORefut}{\operatorname{Refut}}
\providecommand{\ORProv}{\operatorname{Prov}}
\providecommand{\cutrank}{\operatorname{cr}}
\providecommand{\Refut}{\operatorname{Refut}}
\providecommand{\DischargeRule}{\mathrm{Dis}}
\providecommand{\Contraction}{\mathrm{C}}
\providecommand{\bottomAlignProof}{\relax}
\providecommand{\FalseCl}{\bot}
\providecommand{\red}{\mathrel{\rightarrow}}
\providecommand{\varAssign}{\sigma}
\providecommand{\defiff}{\mathrel{\leftrightarrow}}
\providecommand{\FalseInt}{\bot}
\providecommand{\fn}{\operatorname{fn}}
\providecommand{\dotted}{\cdot}

% --- tag-macro semantic fix (2026-10-13, final): \iftag CHOOSES a branch and the OLP proof body is
% always the SECOND group (\iftag{tag}{body}, {tag}{body}{c}, {tag}{body}{c}{d}); a definition that
% prints the tag literal drops the proof body. Bodies contain blank lines, hence \long.
\makeatletter
\providecommand{\iftag}{}
\long\def\iftag#1#2{\@ifnextchar\bgroup{\iftag@three{#2}}{\iftag@stop{#2}}}
\long\def\iftag@three#1#2{\iftag@stop{#1}\@ifnextchar\bgroup{\iftag@eat}{\relax}}
\long\def\iftag@eat#1{}
\long\def\iftag@stop#1{#1}
\long\def\oliflabeldef#1#2{\@ifnextchar\bgroup{\@gobble}{\iftag@stop{#2}}}
\providecommand{\tagrefs}[1]{}
\providecommand{\settag}[1]{}
\makeatother
\begin{document}

% ==== OLP file: content/normal-modal-logic/tableaux/completeness.tex ====
% Part: normal-modal-logic
% Chapter: tableaux
% Section: completeness



\olfileid{nml}{tab}{cpl}

\olsection{Completeness for \Log{K}}

\begin{explain}
  To show that the method of \olup{tableau}s is complete, we have to show
  that whenever there is no closed \olup{tableau} to show $\Gamma \Proves
  A$, then $\Gamma \Entails/ A$, i.e., there is a countermodel. But
  ``there is no closed \olup{tableau}'' means that every way we could try
  to construct one has to fail to close. The trick is to see that if
  every such way fails to close, then a specific, \emph{systematic and
    exhaustive} way also fails to close. And this systematic and
  exhaustive way would close if a closed \olup{tableau} exists. The
  single tableau will contain, among its open branches, all the
  information required to define a countermodel. The countermodel
  given by an open branch in this tableau will contain the all the
  prefixes used on that branch as the worlds, and \olup{propositional
    variable}~$p$ is true at $\sigma$ iff $\sFmla{\True}{p}[\sigma]$
  occurs on the branch.
\end{explain}

\begin{defn}
  A branch in \olup{tableau} is called complete if, whenever it contains
  a prefixed \olup{formula} $\sFmla{S}{A}[\sigma]$ to which a rule can
  be applied, it also contains
  \begin{enumerate}
    \item the prefixed \olup{formula}s that are the corresponding
      conclusions of the rule, in the case of propositional stacking
      rules;
    \item one of the corresponding conclusion \olup{formula}s in the case
      of propositional branching rules;
    \item at least one possible conclusion in the case of modal rules
      that require a new prefix;
    \item the corresponding conclusion for every prefix occurring on
      the branch in the case of modal rules that require a used
      prefix.
  \end{enumerate}
\end{defn}

\begin{explain}
For instance, a complete branch contains $\sFmla{\True}{B}[\sigma]$
and $\sFmla{\True}{C}[\sigma]$ whenever it contains $\sFmla{\True}{B
  \land C}$. If it contains $\sFmla{\True}{B \lor C}[\sigma]$ it
contains at least one of $\sFmla{\False}{B}[\sigma]$ and
$\sFmla{\True}{C}[\sigma]$. If it contains \iftag{prvBox}
{$\sFmla{\False}{\Box}[\sigma]$} {$\sFmla{\True}{\Diamond}[\sigma]$}
it also contains \iftag{prvBox} {$\sFmla{\False}{\Box}[\sigma.n]$}
{$\sFmla{\True}{\Diamond}[\sigma.n]$} for at least one~$n$.  And
whenever it contains \iftag{prvBox} {$\sFmla{\True}{\Box}[\sigma]$}
{$\sFmla{\False}{\Diamond}[\sigma]$} it also contains \iftag{prvBox}
{$\sFmla{\True}{\Box}[\sigma.n]$}
{$\sFmla{\False}{\Diamond}[\sigma.n]$} for every~$n$ such that
$\sigma.n$ is used on the branch.
\end{explain}

\begin{prop}\ollabel{prop:complete-tableau}
  Every finite $\Gamma$ has \olup{tableau} in which every branch is complete.
\end{prop}

\begin{proof}
  Consider an open branch in \olup{tableau} for~$\Gamma$. There are
  finitely many prefixed \olup{formula}s in the branch to which a rule
  could be applied. In some fixed order (say, top to bottom), for each
  of these prefixed \olup{formula}s for which the conditions (1)--(4) do
  not already hold, apply the rules that can be applied to it to
  extend the branch. In some cases this will result in branching;
  apply the rule at the tip of each resulting branch for all remaining
  prefixed \olup{formula}s. Since the number of prefixed \olup{formula}s is
  finite, and the number of used prefixes on the branch is finite,
  this procedure eventually results in (possibly many) branches
  extending the original branch. Apply the procedure to each, and
  repeat. But by construction, every branch is closed.
\end{proof}

\begin{thm}[Completeness]
  \ollabel{thm:tableau-completeness}
  If $\Gamma$ has no closed \olup{tableau}, $\Gamma$ is satisfiable.
\end{thm}

\begin{proof}
By the proposition, $\Gamma$ has \olup{tableau} in which every branch is
complete. Since it has no closed \olup{tableau}, it thas has \olup{tableau} in
which at least one branch is open and complete. Let $\Delta$ be the
set of prefixed \olup{formula}s on the branch, and $P(\Delta)$ the set of
prefixes occurring in it.

We define a model~$\mModel{M(\Delta)} = \tuple{P(\Delta), R, V}$ where
the worlds are the prefixes occurring in~$\Delta$, the accessibility
relation is given by:
\[
R\sigma\sigma' \quad \text{iff} \quad
\sigma'=\sigma.n \quad \text{for some~$n$}
\]
and
\[
V(p) = \Setabs{\sigma}{\sFmla{\True}{p}[\sigma] \in \Delta}.
\]
We show by induction on~$A$ that if $\sFmla{\True}{A}[\sigma] \in
\Delta$ then $\mSat{M(\Delta)}{A}[\sigma]$, and if
$\sFmla{\False}{A}[\sigma] \in \Delta$ then
$\mSat/{M(\Delta)}{A}[\sigma]$.

\begin{enumerate}
  \item \indcase{A}{p}{If $\sFmla{\True}{\indfrm}[\sigma] \in \Delta$
    then $\sigma \in V(p)$ (by definition of~$V$) and so
    $\mSat{M(\Delta)}{\indfrm}[\sigma]$.

    If $\sFmla{\False}{\indfrm}[\sigma] \in \Delta$ then
    $\sFmla{\True}{\indfrm}[\sigma] \notin \Delta$, since the branch
    would otherwise be closed. So $\sigma \notin V(p)$ and thus
    $\mSat/{M(\Delta)}{\indfrm}[\sigma]$.}
  \item \indcase{A}{\lnot B}{\iftag{probNot}{Exercise.}{If
      $\sFmla{\True}{\indfrm}[\sigma] \in \Delta$, then
      $\sFmla{\False}{B}[\sigma] \in \Delta$ since the branch is
      complete. By induction hypothesis,
      $\mSat/{M(\Delta)}{B}[\sigma]$ and thus
      $\mSat{M(\Delta)}{\indfrm}[\sigma]$.

      If $\sFmla{\False}{\indfrm}[\sigma] \in \Delta$, then
      $\sFmla{\True}{B}[\sigma] \in \Delta$ since the branch is
      complete. By induction hypothesis,
      $\mSat{M(\Delta)}{B}[\sigma]$ and thus
      $\mSat/{M(\Delta)}{\indfrm}[\sigma]$.}}
  \item \indcase{A}{B \land C}{\iftag{probAnd}{Exercise.}{If
      $\sFmla{\True}{\indfrm}[\sigma] \in \Delta$, then both
      $\sFmla{\True}{B}[\sigma] \in \Delta$ and
      $\sFmla{\True}{C}[\sigma] \in \Delta$ since the branch is
      complete. By induction hypothesis,
      $\mSat{M(\Delta)}{B}[\sigma]$ and
      $\mSat{M(\Delta)}{C}[\sigma]$. Thus
      $\mSat{M(\Delta)}{\indfrm}[\sigma]$.

      If $\sFmla{\False}{\indfrm}[\sigma] \in \Delta$, then either
      $\sFmla{\False}{B}[\sigma] \in \Delta$ or
      $\sFmla{\False}{C}[\sigma] \in \Delta$ since the branch is
      complete. By induction hypothesis, either
      $\mSat/{M(\Delta)}{B}[\sigma]$ or
      $\mSat/{M(\Delta)}{B}[\sigma]$. Thus
      $\mSat/{M(\Delta)}{\indfrm}[\sigma]$.}}
  \item \indcase{A}{B \lor C}{\iftag{probOr}{Exercise.}{If
      $\sFmla{\True}{\indfrm}[\sigma] \in \Delta$, then either
      $\sFmla{\True}{B}[\sigma] \in \Delta$ or
      $\sFmla{\True}{C}[\sigma] \in \Delta$ since the branch is
      complete. By induction hypothesis, either
      $\mSat{M(\Delta)}{B}[\sigma]$ or
      $\mSat{M(\Delta)}{C}[\sigma]$. Thus
      $\mSat{M(\Delta)}{\indfrm}[\sigma]$.

      If $\sFmla{\False}{\indfrm}[\sigma] \in \Delta$, then both
      $\sFmla{\False}{B}[\sigma] \in \Delta$ and
      $\sFmla{\False}{C}[\sigma] \in \Delta$ since the branch is
      complete. By induction hypothesis, both
      $\mSat/{M(\Delta)}{B}[\sigma]$ and
      $\mSat/{M(\Delta)}{B}[\sigma]$. Thus
      $\mSat/{M(\Delta)}{\indfrm}[\sigma]$.}}
  \item \indcase{A}{B \lif C}{\iftag{probIf}{Exercise.}{If
      $\sFmla{\True}{\indfrm}[\sigma] \in \Delta$, then either
      $\sFmla{\False}{B}[\sigma] \in \Delta$ or
      $\sFmla{\True}{C}[\sigma] \in \Delta$ since the branch is
      complete. By induction hypothesis, either
      $\mSat/{M(\Delta)}{B}[\sigma]$ or
      $\mSat{M(\Delta)}{C}[\sigma]$. Thus
      $\mSat{M(\Delta)}{\indfrm}[\sigma]$.

      If $\sFmla{\False}{\indfrm}[\sigma] \in \Delta$, then both
      $\sFmla{\True}{B}[\sigma] \in \Delta$ and
      $\sFmla{\False}{C}[\sigma] \in \Delta$ since the branch is
      complete. By induction hypothesis, both
      $\mSat{M(\Delta)}{B}[\sigma]$ and
      $\mSat/{M(\Delta)}{B}[\sigma]$. Thus
      $\mSat/{M(\Delta)}{\indfrm}[\sigma]$.}}
  \item \indcase{A}{\Box B}{\iftag{probBox}{Exercise.}{If
      $\sFmla{\True}{\indfrm}[\sigma] \in \Delta$, then, since the
      branch is complete, $\sFmla{\True}{B}[\sigma.n] \in \Delta$ for
      every $\sigma.n$ used on the branch, i.e., for every $\sigma'
      \in P(\Delta)$ such that $R\sigma\sigma'$.  By induction
      hypothesis, $\mSat{M(\Delta)}{B}[\sigma']$ for every $\sigma'$
      such that $R\sigma\sigma'$. Therefore,
      $\mSat{M(\Delta)}{\indfrm}[\sigma]$.

      If $\sFmla{\False}{\indfrm}[\sigma] \in \Delta$, then for some
      $\sigma.n$, $\sFmla{\False}{B}[\sigma.n] \in \Delta$ since the
      branch is complete.  By induction hypothesis,
      $\mSat/{M(\Delta)}{B}[\sigma.n]$.  Since $R\sigma(\sigma.n)$,
      there is a~$\sigma'$ such that
      $\mSat/{M(\Delta)}{B}[\sigma']$. Thus
      $\mSat/{M(\Delta)}{\indfrm}[\sigma]$.}}
  \item \indcase{A}{\Diamond B}{\iftag{probDiamond}{Exercise.}{If
      $\sFmla{\True}{\indfrm}[\sigma] \in \Delta$, then for some
      $\sigma.n$, $\sFmla{\True}{B}[\sigma.n] \in \Delta$ since the
      branch is complete. By induction hypothesis,
      $\mSat{M(\Delta)}{B}[\sigma.n]$.  Since $R\sigma(\sigma.n)$,
      there is a~$\sigma'$ such that
      $\mSat{M(\Delta)}{B}[\sigma']$. Thus
      $\mSat{M(\Delta)}{\indfrm}[\sigma]$.

      If $\sFmla{\False}{\indfrm}[\sigma] \in \Delta$, then, since the
      branch is complete, $\sFmla{\False}{B}[\sigma.n] \in \Delta$
      for every $\sigma.n$ used on the branch, i.e., for every
      $\sigma' \in P(\Delta)$ such that $R\sigma\sigma'$.  By
      induction hypothesis, $\mSat/{M(\Delta)}{B}[\sigma']$ for every
      $\sigma'$ such that $R\sigma\sigma'$. Therefore,
      $\mSat/{M(\Delta)}{\indfrm}[\sigma]$.}}
\end{enumerate}
Since $\Gamma \subseteq \Delta$, $\mSat{M(\Delta)}{\Gamma}$.
\end{proof}

\begin{prob}
Complete the proof of \olref[nml][tab][cpl]{thm:tableau-completeness}.
\end{prob}

\begin{cor}
\ollabel{cor:entailment-completeness}
If $\Gamma \Entails A$ then $\Gamma \Proves A$.
\end{cor}

\begin{cor}
\ollabel{cor:weak-completeness}
If $A$ is true in all models, then $\Proves A$.
\end{cor}


% ==== OLP file: content/normal-modal-logic/tableaux/countermodels.tex ====
% Part: normal-modal-logic
% Chapter: tableaux
% Section: countermodels



\olfileid{nml}{tab}{cou}
      
\olsection{Countermodels from \usetoken{P}{tableau}}

\begin{explain}
  The proof of the completeness theorem doesn't just show that if
  $\Entails A$ then $\Proves A$, it also gives us a method for
  constructing countermodels to~$A$ if $\Entails/ A$. In the case
  of~$\Log{K}$, this method constitutes a \emph{decision
    procedure}. For suppose $\Entails/ A$. Then the proof of
  \olref[cpl]{prop:complete-tableau} gives a method for constructing a
  complete \olup{tableau}. The method in fact always terminates. The
  propositional rules for~$\Log{K}$ only add prefixed \olup{formula}s of
  lower complexity, i.e., each propositional rule need only be applied
  once on a branch for any signed formula $\sFmla{S}{A}[\sigma]$. New
  prefixes are only generated by the
  \iftag{prvBox}{$\TRule{\False}{\Box}$}{}\iftag{notprvBox,notprvDiamond}{}{
    and }\iftag{prvDiamond}{$\TRule{\True}{\Diamond}$}{}
  \iftag{notprvBox,notprvDiamond}{rule}{rules}, and also only have to
  be applied once (and produce a single new prefix).
  \iftag{prvBox}{$\TRule{\True}{\Box}$}{}\iftag{notprvBox,notprvDiamond}{}{
    and }\iftag{prvDiamond}{$\TRule{\False}{\Diamond}$}{}
  \iftag{notprvBox,notprvDiamond}{has}{have} to be applied potentially
  multiple times, but only once per prefix, and only finitely many new
  prefixes are generated. So the construction either results in a
  closed branch or a complete branch after finitely many stages.

  Once a tableau with an open complete branch is constructed, the
  proof of \olref[cpl]{thm:tableau-completeness} gives us an explict model
  that satisfies the original set of prefixed \olup{formula}s. So not
  only is it the case that if $\Gamma \Entails A$, then a closed
  \olup{tableau} exists and $\Gamma \Proves A$, if we look for the
  closed \olup{tableau} in the right way and end up with a ``complete''
  \olup{tableau}, we'll not only know that $\Gamma \Entails/ A$ but
  actually be able to construct a countermodel.
\end{explain}

\iftag{prvBox}{
\begin{ex}
  We know that $\Proves/ \Box(p \lor q) \lif (\Box p \lor \Box
  q)$. The construction of a tableau begins with:
  \begin{oltableau}
    [\pFmla{\False}{\Box(p \lor q) \lif (\Box p \lor \Box q)}{1},
      just = \TAss, checked
      [\pFmla{\True}{\Box(p \lor q)}{1},
        just = {\TRule{\False}{\lif}[1]}, 
        [\pFmla{\False}{\Box p \lor \Box q}{1},
          just = {\TRule{\False}{\lif}[1]}, checked
          [\pFmla{\False}{\Box p}{1},
            just = {\TRule{\False}{\lor}[3]}, checked
            [\pFmla{\False}{\Box q}{1},
              just = {\TRule{\False}{\lor}[3]}, checked
              [\pFmla{\False}{p}{1.1}, 
                just = {\TRule{\False}{\Box}[4]}, checked
                [\pFmla{\False}{q}{1.2}, 
                  just = {\TRule{\False}{\Box}[5]}, checked
                ]
              ]
            ]
          ]
        ]
      ]
    ]
  \end{oltableau}
  The \olup{tableau} is of course not finished yet. In the next step, we
  consider the only line without a checkmark: the prefixed \olup{formula}
  $\sFmla{\True}{\Box(p \lor q)}[1]$ on line~$2$. The construction of
  the closed tableau says to apply the $\TRule{\True}{\Box}$ rule for
  every prefix used on the branch, i.e., for both $1.1$ and $1.2$:
  \begin{oltableau}
    [\pFmla{\False}{\Box(p \lor q) \lif (\Box p \lor \Box q)}{1},
      just = \TAss, checked
      [\pFmla{\True}{\Box(p \lor q)}{1},
        just = {\TRule{\False}{\lif}[1]}, 
        [\pFmla{\False}{\Box p \lor \Box q}{1},
          just = {\TRule{\False}{\lif}[1]}, checked
          [\pFmla{\False}{\Box p}{1},
            just = {\TRule{\False}{\lor}[3]}, checked
            [\pFmla{\False}{\Box q}{1},
              just = {\TRule{\False}{\lor}[3]}, checked
              [\pFmla{\False}{p}{1.1}, 
                just = {\TRule{\False}{\Box}[4]}, checked
                [\pFmla{\False}{q}{1.2}, 
                  just = {\TRule{\False}{\Box}[5]}, checked
                  [\pFmla{\True}{p \lor q}{1.1}, 
                    just = {\TRule{\True}{\Box}[2]}
                    [\pFmla{\True}{p \lor q}{1.2}, 
                      just = {\TRule{\True}{\Box}[2]}
                    ]
                  ]
                ]
              ]
            ]
          ]
        ]
      ]
    ]
  \end{oltableau}
  Now lines 2, 8, and 9, don't have checkmarks. But no new prefix has
  been added, so we apply $\TRule{\True}{\lor}$ to lines~8 and~9, on
  all resulting branches (as long as they don't close):
  \begin{oltableau}
    [\pFmla{\False}{\Box(p \lor q) \lif (\Box p \lor \Box q)}{1},
      just = \TAss, checked
      [\pFmla{\True}{\Box(p \lor q)}{1},
        just = {\TRule{\False}{\lif}[1]}, checked
        [\pFmla{\False}{\Box p \lor \Box q}{1},
          just = {\TRule{\False}{\lif}[1]}, checked
          [\pFmla{\False}{\Box p}{1},
            just = {\TRule{\False}{\lor}[3]}, checked
            [\pFmla{\False}{\Box q}{1},
              just = {\TRule{\False}{\lor}[3]}, checked
              [\pFmla{\False}{p}{1.1}, 
                just = {\TRule{\False}{\Box}[4]}, checked
                [\pFmla{\False}{q}{1.2}, 
                  just = {\TRule{\False}{\Box}[5]}, checked
                  [\pFmla{\True}{p \lor q}{1.1}, 
                    just = {\TRule{\True}{\Box}[2]}, checked
                    [\pFmla{\True}{p \lor q}{1.2}, 
                      just = {\TRule{\True}{\Box}[2]}, checked
                      [\pFmla{\True}{p}{1.1},
                        just = {\TRule{\True}{\lor}[8]}, checked, close
                      ]
                      [\pFmla{\True}{q}{1.1},
                        just = {\TRule{\True}{\lor}[8]}, checked
                        [\pFmla{\True}{p}{1.2},
                          just = {\TRule{\True}{\lor}[9]}, checked]
                        [\pFmla{\True}{q}{1.2},
                          just = {\TRule{\True}{\lor}[9]}, checked, close]
                      ]
                    ]
                  ]
                ]
              ]
            ]
          ]
        ]
      ]
    ]
  \end{oltableau}
  There is one remaining open branch, and it is complete. From it we
  define the model with worlds $W = \{1, 1.1, 1.2\}$ (the only
  prefixes appearing on the open branch), the accessibility relation
  $R = \{\tuple{1, 1.1}, \tuple{1, 1.2}\}$, and the assignment $V(p) =
  \{1.2\}$ (because line~11 contains $\sFmla{\True}{p}[1.2]$) and
  $V(q) = \{1.1\}$ (because line~10 contains
  $\sFmla{\True}{q}[1.1]$). The model is pictured in
  \olref{fig:counter-Box}, and you can verify that it is a
  countermodel to $\Box(p \lor q) \lif (\Box p \lor \Box q)$.
  \begin{figure}
  \begin{center}
    \begin{tikzpicture}[modal]
      \node[world] (w1) [label={[align=right]right:\mFalse{p}\\ \mFalse{q}}]{$1$}; 
      \node[world] (w2) [label={[align=right]right:\mFalse{p}\\ \mTrue{q}},
        above left=of w1]{$1.1$}; 
      \node[world] (w3) [label={[align=right]right:\mTrue{p}\\ \mFalse{q}},
        above right=of w1] {$1.2$};
      \draw[->] (w1) to (w2);
      \draw[->] (w1) to (w3);
    \end{tikzpicture}
  \end{center}
  \caption{A countermodel to $\Box(p \lor q) \lif (\Box p \lor \Box
    q)$.}
  \ollabel{fig:counter-Box}
\end{figure}
\end{ex}          
}
{
\begin{ex}
  We know that $\Proves/ (\Diamond p \land \Diamond q) \lif \Diamond(p
  \land q)$. The construction of a tableau begins with:
  \begin{oltableau}
    [\pFmla{\False}{(\Diamond p \land \Diamond q) \lif \Diamond(p \land q)}{1},
      just = \TAss, checked
      [\pFmla{\True}{\Diamond p \land \Diamond q}{1},
        just = {\TRule{\False}{\lif}[1]}, checked
        [\pFmla{\False}{\Diamond(p \land q)}{1},
          just = {\TRule{\False}{\lif}[1]}
          [\pFmla{\True}{\Diamond p}{1},
            just = {\TRule{\True}{\land}[2]}, checked
            [\pFmla{\True}{\Diamond q}{1},
              just = {\TRule{\True}{\land}[2]}, checked
              [\pFmla{\True}{p}{1.1}, 
                just = {\TRule{\True}{\Diamond}[4]}, checked
                [\pFmla{\True}{q}{1.2}, 
                  just = {\TRule{\True}{\Diamond}[5]}, checked
                ]
              ]
            ]
          ]
        ]
      ]
    ]
  \end{oltableau}
  The \olup{tableau} is of course not finished yet. In the next step, we
  consider the only line without a checkmark: the prefixed \olup{formula}
  $\sFmla{\True}{\Diamond(p \land q)}[1]$ on line~$3$. The construction of
  the closed tableau says to apply the $\TRule{\True}{\Diamond}$ rule for
  every prefix used on the branch, i.e., for both $1.1$ and $1.2$:
  \begin{oltableau}
    [\pFmla{\False}{\Diamond(p \land q) \lif (\Diamond p \land \Diamond q)}{1},
      just = \TAss, checked
      [\pFmla{\True}{\Diamond p \land \Diamond q}{1},
        just = {\TRule{\False}{\lif}[1]}, checked
        [\pFmla{\False}{\Diamond (p \land q)}{1},
          just = {\TRule{\False}{\lif}[1]}
          [\pFmla{\True}{\Diamond p}{1},
            just = {\TRule{\True}{\land}[2]}, checked
            [\pFmla{\True}{\Diamond q}{1},
              just = {\TRule{\True}{\land}[2]}, checked
              [\pFmla{\True}{p}{1.1}, 
                just = {\TRule{\True}{\Diamond}[4]}, checked
                [\pFmla{\True}{q}{1.2}, 
                  just = {\TRule{\True}{\Diamond}[5]}, checked
                  [\pFmla{\False}{p \land q}{1.1}, 
                    just = {\TRule{\False}{\Diamond}[3]}
                    [\pFmla{\False}{p \land q}{1.2}, 
                      just = {\TRule{\False}{\Diamond}[3]}
                    ]
                  ]
                ]
              ]
            ]
          ]
        ]
      ]
    ]
  \end{oltableau}
  Now lines 3, 8, and 9, don't have checkmarks. But no new prefix has
  been added, so we apply $\TRule{\False}{\land}$ to lines~8 and~9, on
  all resulting branches (as long as they don't close):
  \begin{oltableau}
    [\pFmla{\False}{(\Diamond p \land \Diamond q) \lif \Diamond(p \land q)}{1},
      just = \TAss, checked
      [\pFmla{\True}{\Diamond p \land \Diamond q}{1},
        just = {\TRule{\False}{\lif}[1]}, checked
        [\pFmla{\False}{\Diamond(p \land q)}{1},
          just = {\TRule{\False}{\lif}[1]}
          [\pFmla{\True}{\Diamond p}{1},
            just = {\TRule{\True}{\land}[2]}, checked
            [\pFmla{\True}{\Diamond q}{1},
              just = {\TRule{\True}{\land}[2]}, checked
              [\pFmla{\True}{p}{1.1}, 
                just = {\TRule{\True}{\Diamond}[4]}, checked
                [\pFmla{\True}{q}{1.2}, 
                  just = {\TRule{\True}{\Diamond}[5]}, checked
                  [\pFmla{\False}{p \land q}{1.1}, 
                    just = {\TRule{\False}{\Diamond}[3]}, checked
                    [\pFmla{\False}{p \land q}{1.2}, 
                      just = {\TRule{\False}{\Diamond}[3]}, checked
                      [\pFmla{\False}{p}{1.1},
                        just = {\TRule{\False}{\land}[8]}, checked, close
                      ]
                      [\pFmla{\False}{q}{1.1},
                        just = {\TRule{\False}{\land}[8]}, checked
                        [\pFmla{\False}{p}{1.2},
                          just = {\TRule{\False}{\land}[9]}, checked]
                        [\pFmla{\False}{q}{1.2},
                          just = {\TRule{\False}{\land}[9]}, checked, close]
                      ]
                    ]
                  ]
                ]
              ]
            ]
          ]
        ]
      ]
    ]
  \end{oltableau}
  There is one remaining open branch, and it is complete. From it we
  define the model with worlds $W = \{1, 1.1, 1.2\}$ (the only
  prefixes appearing on the open branch), the accessibility relation
  $R = \{\tuple{1, 1.1}, \tuple{1, 1.2}\}$, and the assignment $V(p) =
  \{1.1\}$ (because line~6 contains $\sFmla{\True}{p}[1.1]$) and
  $V(q) = \{1.2\}$ (because line~7 contains
  $\sFmla{\True}{q}[1.1]$). The model is pictured in
  \olref{fig:counter-Diamond}, and you can verify that it is a
  countermodel to $(\Diamond p \land \Diamond q) \lif \Diamond (p \land q)$.
  \begin{figure}
  \begin{center}
    \begin{tikzpicture}[modal]
      \node[world] (w1) [label={[align=right]right:\mFalse{p}\\ \mFalse{q}}]{$1$}; 
      \node[world] (w2) [label={[align=right]right:\mTrue{p}\\ \mFalse{q}},
        above left=of w1]{$1.1$}; 
      \node[world] (w3) [label={[align=right]right:\mFalse{p}\\ \mTrue{q}},
        above right=of w1] {$1.2$};
      \draw[->] (w1) to (w2);
      \draw[->] (w1) to (w3);
    \end{tikzpicture}
  \end{center}
  \caption{A countermodel to $(\Diamond p \land \Diamond q) \lif
    \Diamond (p \land q)$.}
  \ollabel{fig:counter-Diamond}
\end{figure}
\end{ex}          
}


% ==== OLP file: content/normal-modal-logic/tableaux/introduction.tex ====
% Part: normal-modal-logic
% Chapter: tableaux
% Section: introduction



\olfileid{nml}{tab}{int}

\olsection{Introduction}

!!^{tableau}s are certain (downward-branching) trees of \olup{signed
  formula}s, i.e., pairs consisting of a truth value sign ($\True$ or
$\False$) and \olup{sentence}
\[
\sFmla{\True}{A} \text{ or } \sFmla{\False}{A}.
\]
!!^a{tableau} begins with a number of \emph{assumptions}. Each further
\olup{signed formula} is generated by applying one of the inference
rules. Some inference rules add one or more \olup{signed formula}s to a
tip of the tree; others add two new tips, resulting in two branches.
Rules result in \olup{signed formula}s where the \olup{formula} is
less complex than that of the \olup{signed formula} to which it was
applied. When a branch contains both $\sFmla{\True}{A}$ and
$\sFmla{\False}{A}$, we say the branch is \emph{closed}. If every
branch in \olup{tableau} is closed, the entire \olup{tableau} is closed. A
closed \olup{tableau} constitutes \olup{derivation} that shows that the set
of \olup{signed formula}s which were used to begin the \olup{tableau} are
unsatisfiable.  This can be used to define a $\Proves$ relation:
$\Gamma \Proves A$ iff there is some finite set~$\Gamma_0 = \{B_1,
\dots, B_n\} \subseteq \Gamma$ such that there is a closed
\olup{tableau} for the assumptions
\[
\{\sFmla{\False}{A}, \sFmla{\True}{B_1}, \dots, \sFmla{\True}{B_n}\}.
\]

For modal logics, we have to both extend the notion of \olup{signed
formula} and add rules that
cover~\iftag{prvBox}{$\Box$\iftag{prvDiamond}{ and
    $\Diamond$}}{$\Diamond$}. In addition to a sign($\True$ or
$\False$), \olup{formula}s in modal \olup{tableau}s also have
\emph{prefixes}~$\sigma$. The prefixes are non-empty sequences of
positive integers, i.e., $\sigma \in (\PosInt)^* \setminus
\{\emptyseq\}$. When we write such prefixes without the surrounding
$\tuple{\ }$, and separate the individual \olup{element}s by~$.$'s
instead of $,$'s. If $\sigma$ is a prefix, then $\sigma.n$ is $\sigma
\concat \tuple{n}$; e.g., if $\sigma = 1.2.1$, then $\sigma.3$ is
$1.2.1.3$. So for instance,
\[
\sFmla{\True}{\Box A \lif A}[1.2]
\]
is a \emph{prefixed \olup{signed formula}} (or just a \emph{prefixed
  \olup{formula}} for short).

Intuitively, the prefix names a world in a model that might satisfy
the \olup{formula}s on a branch of \olup{tableau}, and if $\sigma$ names
some world, then $\sigma.n$ names a world accessible from (the world
named by)~$\sigma$.


% ==== OLP file: content/normal-modal-logic/tableaux/more-rules.tex ====
% Part: normal-modal-logic
% Chapter: tableaux
% Section: more-rules



\olfileid{nml}{tab}{mru}

\olsection{Rules for Other Accessibility Relations}

In order to deal with logics determined by special accessibility
relations, we consider the additional rules in \olref{tab:more-rules}.

\begin{table}
  \begin{center}
    \def\arraystretch{3}\def\fCenter{}
    \iftag{notprvBox,notprvDiamond}
          {\begin{tabular}{|c|}}
          {\begin{tabular}{|c|c|}}
    \hline
    \iftag{prvBox}{
      \AxiomC{\sFmla{\True}{\Box A}[\sigma]}
      \RightLabel{T$\Box$}
      \UnaryInfC{\sFmla{\True}{A}[\sigma]}
      \DisplayProof}{}
    \iftag{notprvBox,notprvDiamond}{}{&}
    \iftag{prvDiamond}{
      \AxiomC{\sFmla{\False}{\Diamond A}[\sigma]}
      \RightLabel{T$\Diamond$}
      \UnaryInfC{\sFmla{\False}{A}[\sigma]}
      \DisplayProof}{}
    \\[1ex]
    \hline
    \iftag{prvBox}{
      \AxiomC{\sFmla{\True}{\Box A}[\sigma]}
      \RightLabel{D$\Box$}
      \UnaryInfC{\sFmla{\True}{\Diamond A}[\sigma]}
      \DisplayProof}{}
    \iftag{notprvBox,notprvDiamond}{}{&}
    \iftag{prvDiamond}{
      \AxiomC{\sFmla{\False}{\Diamond A}[\sigma]}
      \RightLabel{D$\Diamond$}
      \UnaryInfC{\sFmla{\False}{\Box A}[\sigma]}
      \DisplayProof}{}
    \\[1ex]
    \hline
    \iftag{prvBox}{
      \AxiomC{\sFmla{\True}{\Box A}[\sigma.n]}
      \RightLabel{B$\Box$}
      \UnaryInfC{\sFmla{\True}{A}[\sigma]}
      \DisplayProof}{}
    \iftag{notprvBox,notprvDiamond}{}{&}
    \iftag{prvDiamond}{
      \AxiomC{\sFmla{\False}{\Diamond A}[\sigma.n]}
      \RightLabel{B$\Diamond$}
      \UnaryInfC{\sFmla{\False}{A}[\sigma]}
      \DisplayProof}{}
    \\[1ex]
    \hline
    \iftag{prvBox}{
      \AxiomC{\sFmla{\True}{\Box A}[\sigma]}
      \RightLabel{4$\Box$}
      \UnaryInfC{\sFmla{\True}{\Box A}[\sigma.n]}
      \DisplayProof}{}
    \iftag{notprvBox,notprvDiamond}{}{&}
    \iftag{prvDiamond}{
      \AxiomC{\sFmla{\False}{\Diamond A}[\sigma]}
      \RightLabel{4$\Diamond$}
      \UnaryInfC{\sFmla{\False}{\Diamond A}[\sigma.n]}
      \DisplayProof}{}
    \\
    \iftag{notprvBox,notprvDiamond}
         {$\sigma.n$ is used}
         {$\sigma.n$ is used & $\sigma.n$ is used}
    \\[1ex]
    \hline
    \iftag{prvBox}{
      \AxiomC{\sFmla{\True}{\Box A}[\sigma.n]}
      \RightLabel{4r$\Box$}
      \UnaryInfC{\sFmla{\True}{\Box A}[\sigma]}
      \DisplayProof}{}
    \iftag{notprvBox,notprvDiamond}{}{&}
    \iftag{prvDiamond}{
      \AxiomC{\sFmla{\False}{\Diamond A}[\sigma.n]}
      \RightLabel{4r$\Diamond$}
      \UnaryInfC{\sFmla{\False}{\Diamond A}[\sigma]}
      \DisplayProof}{}
    \\[1ex]
    \hline
    \end{tabular}
  \end{center}
  \caption{More modal rules.}
  \ollabel{tab:more-rules}
\end{table}

Adding these rules results in systems that are sound and complete for
the logics given in \olref{tab:logics-rules}.

\begin{table}
  \begin{center}
    \begin{tabular}{lll}
      \hline
      Logic & $R$ is \dots & Rules\\
      \hline
      $\Log{T} = \Log{KT}$ & reflexive &
      \iftag{prvBox}{T$\Box$}{}%
      \iftag{notprvBox,notprvDiamond}{}{, }%
      \iftag{prvDiamond}{T$\Diamond$}{}
      \\ \hline
      $\Log{D} = \Log{KD}$ & serial &
      \iftag{prvBox}{D$\Box$}{}%
      \iftag{notprvBox,notprvDiamond}{}{, }%
      \iftag{prvDiamond}{D$\Diamond$}{}
      \\ \hline
      $\Log{K4}$ & transitive &
      \iftag{prvBox}{4$\Box$}{}%
      \iftag{notprvBox,notprvDiamond}{}{, }%
      \iftag{prvDiamond}{4$\Diamond$}{}
      \\ \hline
      $\Log{B} = \Log{KTB}$ & reflexive, &
      \iftag{prvBox}{T$\Box$}{}%
      \iftag{notprvBox,notprvDiamond}{}{, }%
      \iftag{prvDiamond}{T$\Diamond$}{}\\
      & symmetric &
      \iftag{prvBox}{B$\Box$}{}%
      \iftag{notprvBox,notprvDiamond}{}{, }%
      \iftag{prvDiamond}{B$\Diamond$}{}
      \\ \hline
      $\Log{S4} = \Log{KT4}$ & reflexive, &
      \iftag{prvBox}{T$\Box$}{}%
      \iftag{notprvBox,notprvDiamond}{}{, }%
      \iftag{prvDiamond}{T$\Diamond$}{},\\
      & transitive &
      \iftag{prvBox}{4$\Box$}{}%
      \iftag{notprvBox,notprvDiamond}{}{, }%
      \iftag{prvDiamond}{4$\Diamond$}{}
      \\ \hline
      $\Log{S5} = \Log{KT4B}$ & reflexive, &
      \iftag{prvBox}{T$\Box$}{}%
      \iftag{notprvBox,notprvDiamond}{}{, }%
      \iftag{prvDiamond}{T$\Diamond$}{},\\
      & transitive, &
      \iftag{prvBox}{4$\Box$}{}%
      \iftag{notprvBox,notprvDiamond}{}{, }%
      \iftag{prvDiamond}{4$\Diamond$}{},\\
      &  euclidean &
      \iftag{prvBox}{4r$\Box$}{}%
      \iftag{notprvBox,notprvDiamond}{}{, }%
      \iftag{prvDiamond}{4r$\Diamond$}{}
      \\ \hline
    \end{tabular}
  \end{center}
  \caption{!!^{tableau} rules for various modal logics.}
  \ollabel{tab:logics-rules}
\end{table}


\begin{ex}
  We give a closed tableau that shows $\Log{S5} \Proves \Ax{5}$, i.e.,
  $\Box A \lif \Box\Diamond A$.
  \begin{oltableau}
    [\pFmla{\False}{\Box\formula{A} \lif \Box\Diamond \formula{A}}{1},
      just = \TAss
      [\pFmla{\True}{\Box \formula{A}}{1}, just = {\TRule{\False}{\lif}[1]}
        [\pFmla{\False}{\Box\Diamond \formula{A}}{1},
          just = {\TRule{\False}{\lif}[1]}
          [\pFmla{\False}{\Diamond \formula{A}}{1.1},
            just = {\TRule{\False}{\Box}[3]}
            [\pFmla{\False}{\Diamond \formula{A}}{1},
              just = {4r$\Diamond$ 4}
              [\pFmla{\False}{\formula{A}}{1.1},
                just = {\TRule{\False}{\Diamond}[5]}
                [\pFmla{\True}{\formula{A}}{1.1},
                  just = {\TRule{\True}{\Box}[2]}, close]
              ]
            ]
          ]
        ]
      ]
    ]
  \end{oltableau}
\end{ex}

\begin{prob}
Give closed \olup{tableau}s that show the following:
  \begin{enumerate}
  \item $\Log{KT5} \Proves \Ax{B}$;
  \item $\Log{KT5} \Proves \Ax{4}$;
  \item $\Log{KDB4} \Proves \Ax{T}$;
  \item $\Log{KB4} \Proves \Ax{5}$;
  \item $\Log{KB5} \Proves \Ax{4}$;
  \item $\Log{KT} \Proves \Ax{D}$.
  \end{enumerate}
\end{prob}


% ==== OLP file: content/normal-modal-logic/tableaux/more-soundness.tex ====
% Part: normal-modal-logic
% Chapter: tableaux
% Section: more-soundness



\olfileid{nml}{tab}{msn}
      
\olsection{Soundness for Additional Rules}

We say a rule is sound for a class of models if, whenever a branch in
\olup{tableau} is satisfiable in a model from that class, the branch
resulting from applying the rule is also satisfiable in a model from
that class.

\begin{prop}
  \ollabel{prop:soundness-T} 
  \iftag{prvBox}{\Ax{T}$\Box$}{}%
  \iftag{notprvBox,notprvDiamond}{}{ and }%
  \iftag{prvDiamond}{\Ax{T}$\Diamond$}{}
  \iftag{notprvBox,notprvDiamond}{is}{are}
  sound for reflexive models.
\end{prop}

\begin{proof}
\begin{tagenumerate}{prvBox,prvDiamond}
\tagitem{prvBox}{The branch is expanded by applying T$\Box$ to
  $\sFmla{\True}{\Box B}[\sigma] \in \Gamma$:
  \iftag{probBox}{Exercise.}{This results in a new \olup{signed
      formula}~$\sFmla{\True}{B}[\sigma]$ on the branch. Suppose
    $\mSat{M}{\Gamma}[f]$, in particular, $\mSat{M}{\Box
      B}[f(\sigma)]$. Since $R$ is reflexive, we know that
    $Rf(\sigma)f(\sigma)$.  Hence, $\mSat{M}{B}[f(\sigma)]$, i.e.,
    $\mModel{M}, f$ satisfies $\sFmla{\True}{B}[\sigma]$.}}{}
\tagitem{prvDiamond}{The branch is expanded by applying T$\Diamond$ to
  $\sFmla{\False}{\Diamond B}[\sigma] \in \Gamma$:
  \iftag{probDiamond}{Exercise.}{This results in a new \olup{signed
      formula}~$\sFmla{\False}{B}[\sigma]$ on the branch. Suppose
    $\mSat{M}{\Gamma}[f]$, in particular, $\mSat/{M}{\Diamond
      B}[f(\sigma)]$. Since $R$ is reflexive, we know that
    $Rf(\sigma)f(\sigma)$.  Hence, $\mSat/{M}{B}[f(\sigma)]$, i.e.,
    $\mModel{M}, f$ satisfies $\sFmla{\False}{B}[\sigma]$.}}{}
\end{tagenumerate}
\end{proof}

\tagprob{probBox,probDiamond}
\begin{prob}
  Complete the proof of \olref[nml][tab][msn]{prop:soundness-T}
\end{prob}
\tagendprob

\begin{prop}
  \ollabel{prop:soundness-D} 
  \iftag{prvBox}{\Ax{D}$\Box$}{}%
  \iftag{notprvBox,notprvDiamond}{}{ and }%
  \iftag{prvDiamond}{\Ax{D}$\Diamond$}{}
  \iftag{notprvBox,notprvDiamond}{is}{are}
  sound for serial models.
\end{prop}

\begin{proof}
\begin{tagenumerate}{prvBox,prvDiamond}
\tagitem{prvBox}{The branch is expanded by applying D$\Box$ to
  $\sFmla{\True}{\Box B}[\sigma] \in \Gamma$:
  \iftag{probBox}{Exercise.}{This results in a new \olup{signed
      formula}~$\sFmla{\True}{\Diamond B}[\sigma]$ on the
    branch. Suppose $\mSat{M}{\Gamma}[f]$, in particular,
    $\mSat{M}{\Box B}[f(\sigma)]$. Since $R$ is serial, there is a $w
    \in W$ such that $Rf(\sigma)w$. Then $\mSat{M}{B}[w]$, and hence
    $\mSat{M}{\Diamond B}[f(\sigma)]$.  So, $\mModel{M}, f$
    satisfies $\sFmla{\True}{\Diamond B}[\sigma]$.}}{}
\tagitem{prvDiamond}{The branch is expanded by applying D$\Diamond$ to
  $\sFmla{\False}{\Diamond B}[\sigma] \in \Gamma$:
  \iftag{probDiamond}{Exercise.}{This results in a new \olup{signed
      formula}~$\sFmla{\False}{\Box B}[\sigma]$ on the
    branch. Suppose $\mSat{M}{\Gamma}[f]$, in particular,
    $\mSat/{M}{\Diamond B}[f(\sigma)]$. Since $R$ is serial, there is a $w
    \in W$ such that $Rf(\sigma)w$. Then $\mSat/{M}{B}[w]$, and hence
    $\mSat/{M}{\Box B}[f(\sigma)]$.  So, $\mModel{M}, f$
    satisfies $\sFmla{\False}{\Box B}[\sigma]$.}}{}
\end{tagenumerate}
\end{proof}

\tagprob{probBox,probDiamond}
\begin{prob}
  Complete the proof of \olref[nml][tab][msn]{prop:soundness-D}
\end{prob}
\tagendprob

\begin{prop}
  \ollabel{prop:soundness-B} 
  \iftag{prvBox}{\Ax{B}$\Box$}{}%
  \iftag{notprvBox,notprvDiamond}{}{ and }%
  \iftag{prvDiamond}{\Ax{B}$\Diamond$}{}
  \iftag{notprvBox,notprvDiamond}{is}{are}
  sound for symmetric models.
\end{prop}

\begin{proof}
\begin{tagenumerate}{prvBox,prvDiamond}
\tagitem{prvBox}{The branch is expanded by applying B$\Box$ to
  $\sFmla{\True}{\Box B}[\sigma.n] \in \Gamma$:
  \iftag{probBox}{Exercise.}{This results in a new \olup{signed
      formula}~$\sFmla{\True}{B}[\sigma]$ on the branch. Suppose
    $\mSat{M}{\Gamma}[f]$, in particular, $\mSat{M}{\Box
      B}[f(\sigma.n)]$. Since $f$ is an interpretation of prefixes on
    the branch into~$\mModel{M}$, we know that
    $Rf(\sigma)f(\sigma.n)$.  Since $R$ is symmetric,
    $Rf(\sigma.n)f(\sigma)$.  Since $\mSat{M}{\Box B}[f(\sigma.n)]$,
    $\mSat{M}{B}[f(\sigma)]$. Hence, $\mModel{M}, f$ satisfies
    $\sFmla{\True}{B}[\sigma]$.}}{}
\tagitem{prvDiamond}{The branch is expanded by applying B$\Diamond$ to
  $\sFmla{\False}{\Diamond B}[\sigma.n] \in \Gamma$:
  \iftag{probDiamond}{Exercise.}{This results in a new \olup{signed
      formula}~$\sFmla{\False}{B}[\sigma]$ on the branch. Suppose
    $\mSat{M}{\Gamma}[f]$, in particular, $\mSat/{M}{\Diamond
      B}[f(\sigma.n)]$. Since $f$ is an interpretation of prefixes on
    the branch into~$\mModel{M}$, we know that
    $Rf(\sigma)f(\sigma.n)$.  Since $R$ is symmetric,
    $Rf(\sigma.n)f(\sigma)$.  Since $\mSat/{M}{\Diamond B}[f(\sigma.n)]$,
    $\mSat/{M}{B}[f(\sigma)]$. Hence, $\mModel{M}, f$ satisfies
    $\sFmla{\False}{B}[\sigma]$.}}{}
\end{tagenumerate}
\end{proof}

\tagprob{probBox,probDiamond}
\begin{prob}
  Complete the proof of \olref[nml][tab][msn]{prop:soundness-B}
\end{prob}
\tagendprob

\begin{prop}
  \ollabel{prop:soundness-4} 
  \iftag{prvBox}{\Ax{4}$\Box$}{}%
  \iftag{notprvBox,notprvDiamond}{}{ and }%
  \iftag{prvDiamond}{\Ax{4}$\Diamond$}{}
  \iftag{notprvBox,notprvDiamond}{is}{are}
  sound for transitive models.
\end{prop}

\begin{proof}
\begin{tagenumerate}{prvBox,prvDiamond}
\tagitem{prvBox}{The branch is expanded by applying 4$\Box$ to
  $\sFmla{\True}{\Box B}[\sigma] \in \Gamma$:
  \iftag{probBox}{Exercise.}{This results in a new \olup{signed
      formula}~$\sFmla{\True}{\Box B}[\sigma.n]$ on the
    branch. Suppose $\mSat{M}{\Gamma}[f]$, in particular,
    $\mSat{M}{\Box B}[f(\sigma)]$. Since $f$ is an interpretation of
    prefixes on the branch into~$\mModel{M}$ and $\sigma.n$ must be
    used, we know that $Rf(\sigma)f(\sigma.n)$. Now let $w$ be any
    world such that $Rf(\sigma.n)w$.  Since $R$ is transitive,
    $Rf(\sigma)w$.  Since $\mSat{M}{\Box B}[f(\sigma)]$,
    $\mSat{M}{B}[w]$. Hence, $\mSat{M}{\Box B}[f(\sigma.n)]$, and
    $\mModel{M}, f$ satisfies $\sFmla{\True}{\Box B}[\sigma.n]$.}}{}
\tagitem{prvDiamond}{The branch is expanded by applying 4$\Diamond$ to
  $\sFmla{\False}{\Diamond B}[\sigma] \in \Gamma$:
  \iftag{probDiamond}{Exercise.}{This results in a new \olup{signed
      formula}~$\sFmla{\False}{\Diamond B}[\sigma.n]$ on the
    branch. Suppose $\mSat{M}{\Gamma}[f]$, in particular,
    $\mSat/{M}{\Diamond B}[f(\sigma)]$. Since $f$ is an interpretation of
    prefixes on the branch into~$\mModel{M}$ and $\sigma.n$ must be
    used, we know that $Rf(\sigma)f(\sigma.n)$. Now let $w$ be any
    world such that $Rf(\sigma.n)w$.  Since $R$ is transitive,
    $Rf(\sigma)w$.  Since $\mSat/{M}{\Diamond B}[f(\sigma)]$,
    $\mSat/{M}{B}[w]$. Hence, $\mSat/{M}{\Diamond B}[f(\sigma.n)]$, and
    $\mModel{M}, f$ satisfies $\sFmla{\False}{\Diamond B}[\sigma.n]$.}}{}
\end{tagenumerate}
\end{proof}

\tagprob{probBox,probDiamond}
\begin{prob}
  Complete the proof of \olref[nml][tab][msn]{prop:soundness-4}
\end{prob}
\tagendprob

\begin{prop}
  \ollabel{prop:soundness-4r} 
  \iftag{prvBox}{\Ax{4r}$\Box$}{}%
  \iftag{notprvBox,notprvDiamond}{}{ and }%
  \iftag{prvDiamond}{\Ax{4r}$\Diamond$}{}
  \iftag{notprvBox,notprvDiamond}{is}{are}
  sound for euclidean models.
\end{prop}

\begin{proof}
\begin{tagenumerate}{prvBox,prvDiamond}
\tagitem{prvBox}{The branch is expanded by applying 4r$\Box$ to
  $\sFmla{\True}{\Box B}[\sigma.n] \in \Gamma$:
  \iftag{probBox}{Exercise.}{This results in a new \olup{signed
      formula}~$\sFmla{\True}{\Box B}[\sigma]$ on the branch. Suppose
    $\mSat{M}{\Gamma}[f]$, in particular, $\mSat{M}{\Box
      B}[f(\sigma.n)]$. Since $f$ is an interpretation of prefixes on
    the branch into~$\mModel{M}$, we know that
    $Rf(\sigma)f(\sigma.n)$. Now let $w$ be any world such that
    $Rf(\sigma)w$.  Since $R$ is euclidean, $Rf(\sigma.n)w$.  Since
    $\mSat{M}{\Box B}[f(\sigma).n]$, $\mSat{M}{B}[w]$. Hence,
    $\mSat{M}{\Box B}[f(\sigma)]$, and $\mModel{M}, f$ satisfies
    $\sFmla{\True}{\Box B}[\sigma]$.}}{}
\tagitem{prvDiamond}{The branch is expanded by applying 4r$\Diamond$ to
  $\sFmla{\False}{\Diamond B}[\sigma.n] \in \Gamma$:
  \iftag{probDiamond}{Exercise.}{This results in a new \olup{signed
      formula}~$\sFmla{\True}{\Box B}[\sigma]$ on the branch. Suppose
    $\mSat{M}{\Gamma}[f]$, in particular, $\mSat/{M}{\Diamond
      B}[f(\sigma.n)]$. Since $f$ is an interpretation of prefixes on
    the branch into~$\mModel{M}$, we know that
    $Rf(\sigma)f(\sigma.n)$. Now let $w$ be any world such that
    $Rf(\sigma)w$.  Since $R$ is euclidean, $Rf(\sigma.n)w$.  Since
    $\mSat/{M}{\Diamond B}[f(\sigma).n]$, $\mSat/{M}{B}[w]$. Hence,
    $\mSat/{M}{\Diamond B}[f(\sigma)]$, and $\mModel{M}, f$ satisfies
    $\sFmla{\False}{\Diamond B}[\sigma]$.}}{}
\end{tagenumerate}
\end{proof}

\tagprob{probBox,probDiamond}
\begin{prob}
  Complete the proof of \olref[nml][tab][msn]{prop:soundness-4r}
\end{prob}
\tagendprob

\begin{cor}
\ollabel{cor:soundness-logics} The \olup{tableau} systems given in
\olref[mru]{tab:logics-rules} are sound for the respective classes of
models.
\end{cor}


% ==== OLP file: content/normal-modal-logic/tableaux/proofs-in-K.tex ====
% Part: normal-modal-logic
% Chapter: tableaux
% Section: proofs-in-K



\olfileid{nml}{tab}{prk}

\olsection{\usetoken{P}{tableau} for \Log{K}}

\iftag{prvBox}{

\begin{ex}
  We give a closed tableau that shows $\Proves (\Box A \land \Box B)
  \lif \Box (A \land B)$.
  \begin{oltableau}
    [\pFmla{\False}{(\Box\formula{A} \land \Box\formula{B}) \lif
        \Box (\formula{A} \land \formula{B})}{1},
      just =\TAss
      [\pFmla{\True}{\Box\formula{A} \land \Box\formula{B}}{1},
        just = {\TRule{\False}{\lif}[1]}
        [\pFmla{\False}{\Box(\formula{A} \land \formula{B})}{1},
          just = {\TRule{\False}{\lif}[1]}
          [\pFmla{\True}{\Box\formula{A}}{1},
            just = {\TRule{\True}{\land}[2]}
            [\pFmla{\True}{\Box\formula{B}}{1},
              just = {\TRule{\True}{\land}[2]}
              [\pFmla{\False}{\formula{A} \land \formula{B}}{1.1},
                just = {\TRule{\False}{\Box}[3]}
                [\pFmla{\False}{\formula{A}}{1.1},
                  just = {\TRule{\False}{\land}[6]}
                  [\pFmla{\True}{\formula{A}}{1.1},
                    just= {\TRule{\True}{\Box}[4]}, close]]
                [\pFmla{\False}{\formula{B}}{1.1},
                  just = {\TRule{\False}{\land}[6]}
                  [\pFmla{\True}{\formula{B}}{1.1},
                    just= {\TRule{\True}{\Box}[5]}, close]]
              ]
            ]
          ]
        ]
      ]
    ]
  \end{oltableau}
\end{ex}
}{}

\iftag{prvDiamond}{
\begin{ex}
  We give a closed tableau that shows $\Proves \Diamond(A \lor B)
  \lif (\Diamond A \lor \Diamond B)$:
  \begin{oltableau}
    [\pFmla{\False}{\Diamond(\formula{A} \lor \formula{B}) \lif
        (\Diamond \formula{A} \lor \Diamond \formula{B})}{1},
      just =\TAss
      [\pFmla{\True}{\Diamond(\formula{A} \lor \formula{B})}{1},
        just = {\TRule{\False}{\lif}[1]}
        [\pFmla{\False}{\Diamond\formula{A} \lor \Diamond\formula{B}}{1},
          just = {\TRule{\False}{\lif}[1]}
          [\pFmla{\False}{\Diamond\formula{A}}{1},
            just = {\TRule{\False}{\lor}[3]}
            [\pFmla{\False}{\Diamond\formula{B}}{1},
              just = {\TRule{\False}{\lor}[3]}
              [\pFmla{\True}{\formula{A} \lor \formula{B}}{1.1},
                just = {\TRule{\True}{\Diamond}[2]},
                [\pFmla{\True}{\formula{A}}{1.1},
                  just = {\TRule{\True}{\lor}[6]}
                  [\pFmla{\False}{\formula{A}}{1.1},
                    just= {\TRule{\False}{\Diamond}[4]}, close]]
                [\pFmla{\True}{\formula{B}}{1.1},
                  just = {\TRule{\True}{\lor}[6]}
                  [\pFmla{\False}{\formula{B}}{1.1},
                    just= {\TRule{\False}{\Diamond}[5]}, close]]
              ]
            ]
          ]
        ]
      ]
    ]
  \end{oltableau}
\end{ex}
}{}


\begin{prob}
  Find closed \olup{tableau}s in~$\Log{K}$ for the following \olup{formula}s:
  \begin{enumerate}
    \item $\Box \lnot p \lif \Box(p \lif q)$
    \item $(\Box p \lor \Box q) \lif \Box(p \lor q)$
    \item $\Diamond p \lif \Diamond(p \lor q)$
    \item $\Box(p \land q) \lif \Box p$
  \end{enumerate}
\end{prob}


% ==== OLP file: content/normal-modal-logic/tableaux/rules-for-K.tex ====
% Part: normal-modal-logic
% Chapter: tableaux
% Section: rules-for-K



\olfileid{nml}{tab}{rul}

\olsection{Rules for \Ax{K}}

The rules for the regular propositional connectives are the same as
for regular propositional signed \olup{tableau}s, just with prefixes
added. In each case, the rule applied to a signed \olup{formula}
$\sFmla{S}{A}[\sigma]$ produces new \olup{formula}s that are also
prefixed by~$\sigma$. This should be intuitively clear: e.g., if $A
\land B$ is true at (a world named by)~$\sigma$, then $A$ and $B$
are true at~$\sigma$ (and not at any other world). We collect the
propositional rules in \olref{tab:prop-rules}.

\begin{table}
  \[\def\arraystretch{3}\begin{array}{|c|c|}
    \hline
    \AxiomC{\sFmla{\True}{\lnot A}[\sigma]}
    \RightLabel{\TRule{\True}{\lnot}}
    \UnaryInfC{\sFmla{\False}{A}[\sigma]}
    \DisplayProof
    &
    \AxiomC{\sFmla{\False}{\lnot A}[\sigma]}
    \RightLabel{\TRule{\False}{\lnot}}
    \UnaryInfC{\sFmla{\True}{A}[\sigma]}
    \DisplayProof
    \\[1ex]
    \hline
    \AxiomC{\sFmla{\True}{A \land B}[\sigma]}
    \RightLabel{\TRule{\True}{\land}}
    \UnaryInfC{\sFmla{\True}{A}[\sigma]}
    \noLine
    \UnaryInfC{\sFmla{\True}{B}[\sigma]}
    \DisplayProof
    &
    \AxiomC{\sFmla{\False}{A \land B}[\sigma]}
    \RightLabel{\TRule{\False}{\land}}
    \UnaryInfC{$\sFmla{\False}{A}[\sigma] \quad \mid \quad
      \sFmla{\False}{B}[\sigma]$}
    \DisplayProof
    \\[2ex]
    \hline
    \AxiomC{\sFmla{\True}{A \lor B}[\sigma]}
    \RightLabel{\TRule{\True}{\lor}}
    \UnaryInfC{$\sFmla{\True}{A}[\sigma] \quad \mid \quad
      \sFmla{\True}{B}[\sigma]$}
    \DisplayProof
    &
    \AxiomC{\sFmla{\False}{A \lor B}[\sigma]}
    \RightLabel{\TRule{\False}{\lor}}
    \UnaryInfC{\sFmla{\False}{A}[\sigma]}
    \noLine
    \UnaryInfC{\sFmla{\False}{B}[\sigma]}
    \DisplayProof
    \\[2ex]
    \hline
    \AxiomC{\sFmla{\True}{A \lif B}[\sigma]}
    \RightLabel{\TRule{\True}{\lif}}
    \UnaryInfC{$\sFmla{\False}{A}[\sigma] \quad \mid
      \quad \sFmla{\True}{B}[\sigma]$}
    \DisplayProof
    &
    \AxiomC{\sFmla{\False}{A \lif B}[\sigma]}
    \RightLabel{\TRule{\False}{\lif}}
    \UnaryInfC{\sFmla{\True}{A}[\sigma]}
    \noLine
    \UnaryInfC{\sFmla{\False}{B}[\sigma]}
    \DisplayProof
    \\[2ex]
    \hline
  \end{array}\]
  \caption{Prefixed \olup{tableau} rules for the propositional
    connectives}
  \ollabel{tab:prop-rules}
\end{table}

The closure condition is the same as for ordinary \olup{tableau}s,
although we require that not just the \olup{formula}s but also the
prefixes must match. So a branch is closed if it contains both
\[
\sFmla{\True}{A}[\sigma] \quad\text{and}\quad \sFmla{\False}{A}[\sigma]
\]
for some prefix $\sigma$ and \olup{formula}~$A$.

The rules for setting up assumptions is also as for ordinary
\olup{tableau}s, except that for assumptions we always use the
prefix~$1$. (It does not matter which prefix we use, as long as it's
the same for all assumptions.) So, e.g., we say that
\[
B_1, \dots, B_n \Proves A
\]
iff there is a closed tableau for the assumptions
\[
\sFmla{\True}{B_1}[1], \dots, \sFmla{\True}{B_n}[1],
\sFmla{\False}{A}[1].
\]

For the modal operator\iftag{prvBox}{\iftag{prvDiamond}{s~$\Box$
    and}{~$\Box$}}{}\iftag{prvDiamond}{~$\Diamond$}{}, the prefix of
the conclusion of the rule applied to \olup{formula} with
prefix~$\sigma$ is $\sigma.n$. However, which $n$ is allowed depends
on whether the sign is~$\True$ or~$\False$.

\iftag{prvBox}{The $\TRule{\True}{\Box}$ rule extends a branch
  containing $\sFmla{\True}{\Box A}[\sigma]$ by
  $\sFmla{\True}{A}[\sigma.n]$.\iftag{prvDiamond}{ Similarly,
    t}{}}{T}\iftag{prvDiamond}{he $\TRule{\False}{\Diamond}$ rule
  extends a branch containing $\sFmla{\False}{\Diamond A}[\sigma]$ by
  $\sFmla{\False}{A}[\sigma.n]$.}{}
\iftag{notprvBox,notprvDiamond}{It}{They} can only be applied for a
prefix~$\sigma.n$ which \emph{already} occurs on the branch in which
it is applied. Let's call such a prefix ``used'' (on the branch).

\iftag{prvBox}{The $\TRule{\False}{\Box}$ rule extends a branch
  containing $\sFmla{\False}{\Box A}[\sigma]$ by
  $\sFmla{\False}{A}[\sigma.n]$.\iftag{prvDiamond}{ Similarly,
    t}{}}{T}\iftag{prvDiamond}{he $\TRule{\True}{\Diamond}$ rule extends
  a branch containing $\sFmla{\True}{\Diamond A}[\sigma]$ by
  $\sFmla{\True}{A}[\sigma.n]$.}{}
\iftag{notprvBox,notprvDiamond}{This rule}{These rules}, however, can
only be applied for a prefix~$\sigma.n$ which \emph{does not} already
occur on the branch in which it is applied. We call such prefixes
``new'' (to the branch).

The rules are given in \olref{tab:rules-K}.

\begin{table}
  \begin{center}
    \def\arraystretch{3}\def\fCenter{}
    \begin{tabular}{|c|c|}
    \hline
    \iftag{prvBox}{
      \AxiomC{\sFmla{\True}{\Box A}[\sigma]}
      \RightLabel{\TRule{\True}{\Box}}
      \UnaryInfC{\sFmla{\True}{A}[\sigma.n]}
      \DisplayProof
      &
      \AxiomC{\sFmla{\False}{\Box A}[\sigma]}
      \RightLabel{\TRule{\False}{\Box}}
      \UnaryInfC{\sFmla{\False}{A}[\sigma.n]}
      \DisplayProof\\
      $\sigma.n$ is used & $\sigma.n$ is new
      \\[1ex]
      \hline}{}
    \iftag{prvDiamond}{
      \AxiomC{\sFmla{\True}{\Diamond A}[\sigma]}
      \RightLabel{\TRule{\True}{\Diamond}}
      \UnaryInfC{\sFmla{\True}{A}[\sigma.n]}
      \DisplayProof
      &
      \AxiomC{\sFmla{\False}{\Diamond A}[\sigma]}
      \RightLabel{\TRule{\False}{\Diamond}}
      \UnaryInfC{\sFmla{\False}{A}[\sigma.n]}
      \DisplayProof\\
      $\sigma.n$ is new & $\sigma.n$ is used
      \\[1ex]
      \hline}{}
    \end{tabular}
  \end{center}
  \caption{The modal rules for \Ax{K}.}
  \ollabel{tab:rules-K}
\end{table}

The requirement that the restriction that the prefix for
\iftag{prvBox}{\TRule{\True}{\Box}}{\TRule{\False}{\Diamond}} must be
used is necessary as otherwise we would count the following as a
closed \olup{tableau}:
\iftag{notprvBox,notprvDiamond}{%
  \iftag{prvBox}{%
  \begin{oltableau}
    [\pFmla{\True}{\Box \formula{A}}{1}, just = \TAss
      [\pFmla{\False}{\lnot\Box\lnot \formula{A}}{1}, just = \TAss
        [\pFmla{\True}{\formula{A}}{1.1}, just = {\TRule{\True}{\Box}[1]}
          [\pFmla{\True}{\Box\lnot\formula{A}}{1},
            just ={\TRule{\False}{\lnot}[2]}
            [\pFmla{\True}{\lnot\formula{A}}{1.1},
              just = {\TRule{\True}{\Box}[4]}
              [\pFmla{\False}{\formula{A}}{1.1},
                  just ={\TRule{\True}{\lnot}[5]}, close]
            ]
          ]
        ]
      ]
    ]
  \end{oltableau}
  }{
  \begin{oltableau}
    [\pFmla{\True}{\lnot\Diamond\lnot \formula{A}}{1}, just = \TAss
      [\pFmla{\False}{\Diamond\formula{A}}{1}, just = \TAss
        [\pFmla{\False}{\formula{A}}{1.1}, just = {\TRule{\False}{\Diamond}[2]}
          [\pFmla{\False}{\Diamond\lnot\formula{A}}{1},
            just ={\TRule{\True}{\lnot}[1]}
            [\pFmla{\False}{\lnot\formula{A}}{1.1},
              just = {\TRule{\False}{\Diamond}[4]}
              [\pFmla{\True}{\formula{A}}{1.1},
                  just ={\TRule{\True}{\lnot}[5]}, close]
            ]
          ]
        ]
      ]
    ]
  \end{oltableau}
}}{
  \begin{oltableau}
    [\pFmla{\True}{\Box \formula{A}}{1}, just = \TAss
      [\pFmla{\False}{\Diamond \formula{A}}{1}, just = \TAss
        [\pFmla{\True}{\formula{A}}{1.1}, just = {\TRule{\True}{\Box}[1]}
          [\pFmla{\False}{\formula{A}}{1.1},
            just = {\TRule{\False}{\Diamond}[2]}, close]
        ]
      ]
    ]
\end{oltableau}}

But $\Box \formula{A} \Entails/ \Diamond \formula{A}$, so our proof
system would be unsound. Likewise, $\Diamond \formula{A} \Entails/
\Box \formula{A}$, but without the restriction that the prefix for
\iftag{prvBox}{\TRule{\False}{\Box}}{\TRule{\True}{\Diamond}} must be
new, this would be a closed tableau: \iftag{notprvBox,notprvDiamond}{%
  \iftag{prvBox}{%
  \begin{oltableau}
    [\pFmla{\True}{\lnot\Box\lnot \formula{A}}{1}, just = \TAss
      [\pFmla{\False}{\Box \formula{A}}{1}, just = \TAss
        [\pFmla{\False}{\formula{A}}{1.1}, just = {\TRule{\True}{\Box}[2]}
          [\pFmla{\False}{\Box\lnot\formula{A}}{1},
            just ={\TRule{\True}{\lnot}[1]}
            [\pFmla{\False}{\lnot\formula{A}}{1.1},
              just = {\TRule{\False}{\Box}[4]}
              [\pFmla{\True}{\formula{A}}{1.1},
                  just ={\TRule{\False}{\lnot}[5]}, close]
            ]
          ]
        ]
      ]
    ]
  \end{oltableau}
  }{
  \begin{oltableau}
    [\pFmla{\True}{\Diamond \formula{A}}{1}, just = \TAss
      [\pFmla{\False}{\lnot\Diamond\lnot\formula{A}}{1}, just = \TAss
        [\pFmla{\True}{\formula{A}}{1.1}, just = {\TRule{\True}{\Diamond}[1]}
          [\pFmla{\True}{\Diamond\lnot\formula{A}}{1},
            just ={\TRule{\False}{\lnot}[2]}
            [\pFmla{\True}{\lnot\formula{A}}{1.1},
              just = {\TRule{\True}{\Diamond}[4]}
              [\pFmla{\False}{\formula{A}}{1.1},
                  just ={\TRule{\True}{\lnot}[5]}, close]
            ]
          ]
        ]
      ]
    ]
  \end{oltableau}
}}{
  \begin{oltableau}
    [\pFmla{\True}{\Diamond \formula{A}}{1}, just = \TAss,name=one
      [\pFmla{\False}{\Box \formula{A}}{1}, just = \TAss, name=two
        [\pFmla{\True}{\formula{A}}{1.1}, just = {\TRule{\True}{\Diamond}[1]}
          [\pFmla{\False}{\formula{A}}{1.1},
            just = {\TRule{\False}{\Box}[2]}, close]
        ]
      ]
    ]
  \end{oltableau}}


% ==== OLP file: content/normal-modal-logic/tableaux/simple-S5.tex ====
% Part: normal-modal-logic
% Chapter: tableaux
% Section: simple-S5



\olfileid{nml}{tab}{s5}

\olsection{Simple \usetoken{P}{tableau} for \Log{S5}}

\Log{S5} is sound and complete with respect to the class of universal
models, i.e., models where every world is accessible from every
world. In universal models the accessibility relation doesn't matter:
``there is a world~$w$ where $\mSat{M}{A}[w]$'' is true if and only
if there is such a $w$ that's accessible from~$u$. So in \Log{S5}, we
can define models as simply a set of worlds and a valuation~$V$. This
suggests that we should be able to simplify the \olup{tableau} rules as
well. In the general case, we take as prefixes sequences of positive
integers, so that we can keep track of which such prefixes name worlds
which are accessible from others: $\sigma.n$ names a world accessible
from~$\sigma$. But in \Log{S5} any world is accessible from any world,
so there is no need to so keep track. Instead, we can use positive
integers as prefixes. The simplified rules are given in
\olref{tab:rules-S5}.

\begin{table}
  \begin{center}
    \def\arraystretch{3}\def\fCenter{}
    \begin{tabular}{|c|c|}
    \hline
    \iftag{prvBox}{
      \AxiomC{\sFmla{\True}{\Box A}[n]}
      \RightLabel{\TRule{\True}{\Box}}
      \UnaryInfC{\sFmla{\True}{A}[m]}
      \DisplayProof
      &
      \AxiomC{\sFmla{\False}{\Box A}[n]}
      \RightLabel{\TRule{\False}{\Box}}
      \UnaryInfC{\sFmla{\False}{A}[m]}
      \DisplayProof\\
      $m$ is used & $m$ is new
      \\[1ex]
      \hline}{}
    \iftag{prvDiamond}{
      \AxiomC{\sFmla{\True}{\Diamond A}[n]}
      \RightLabel{\TRule{\True}{\Diamond}}
      \UnaryInfC{\sFmla{\True}{A}[m]}
      \DisplayProof
      &
      \AxiomC{\sFmla{\False}{\Diamond A}[n]}
      \RightLabel{\TRule{\False}{\Diamond}}
      \UnaryInfC{\sFmla{\False}{A}[m]}
      \DisplayProof\\
      $m$ is new & $m$ is used
      \\[1ex]
      \hline}{}
    \end{tabular}
  \end{center}
  \caption{Simplified rules for \Log{S5}.}
  \ollabel{tab:rules-S5}
\end{table}

\begin{ex}
  We give a simplified closed tableau that shows $\Log{S5} \Proves
  \Ax{5}$, i.e., $\Diamond A \lif \Box\Diamond A$.
  \begin{oltableau}
    [\pFmla{\False}{\Diamond\formula{A} \lif \Box\Diamond \formula{A}}{1},
      just = \TAss
      [\pFmla{\True}{\Diamond \formula{A}}{1}, just = {\TRule{\False}{\lif}[1]}
        [\pFmla{\False}{\Box\Diamond \formula{A}}{1},
          just = {\TRule{\False}{\lif}[1]}
          [\pFmla{\False}{\Diamond \formula{A}}{2},
            just = {\TRule{\False}{\Box}[3]}
            [\pFmla{\True}{\formula{A}}{3},
              just = {\TRule{\True}{\Diamond}[2]}
                [\pFmla{\False}{\formula{A}}{3},
                  just = {\TRule{\False}{\Diamond}[4]}, close]
            ]
          ]
        ]
      ]
    ]
  \end{oltableau}
\end{ex}


% ==== OLP file: content/normal-modal-logic/tableaux/soundness.tex ====
% Part: normal-modal-logic
% Chapter: tableaux
% Section: soundness



\olfileid{nml}{tab}{sou}

\olsection{Soundness for \Log{K}}

\begin{editorial}
  This soundness proof reuses the soundness proof for classical
  propositional logic, i.e., it proves everything from scratch. That's
  ok if you want a self-contained soundness proof. If you already have
  seen soundness for ordinary tableau this will be repetitive. It's
  planned to make it possible to switch between self-contained version
  and a version building on the non-modal case.
\end{editorial}

\begin{explain}
  In order to show that prefixed \olup{tableau}s are sound, we have to
  show that if
  \[
  \sFmla{\True}{B_1}[1], \dots, \sFmla{\True}{B_n}[1], \sFmla{\False}{A}[1]
  \]
  has a closed \olup{tableau} then $B_1, \dots, B_n \Entails A$. It is
  easier to prove the contrapositive: if for some $\mModel{M}$ and
  world~$w$, $\mSat{M}{B_i}[w]$ for all $i=1$, \dots,~$n$ but
  $\mSat{M}{A}[w]$, then no \olup{tableau} can close. Such a
  countermodel shows that the initial assumptions of the \olup{tableau}
  are satisfiable. The strategy of the proof is to show that whenever
  all the prefixed \olup{formula}s on \olup{tableau} branch are
  satisfiable, any application of a rule results in at least one
  extended branch that is also satisfiable. Since closed branches are
  unsatisfiable, any \olup{tableau} for a satisfiable set of prefixed
  \olup{formula}s must have at least one open branch.

  In order to apply this strategy in the modal case, we have to extend
  our definition of ``satisfiable'' to modal modals and prefixes. With
  that in hand, however, the proof is straightforward.
\end{explain}

\begin{defn}
  Let $P$ be some set of prefixes, i.e., $P \subseteq (\PosInt)^*
  \setminus \{\emptyseq\}$ and let $\mModel{M}$ be a model.  A
  function~$f\colon P \to W$ is an \emph{interpretation of}~$P$
  in~$\mModel{M}$ if, whenever $\sigma$ and $\sigma.n$ are both
  in~$P$, then $Rf(\sigma)f(\sigma.n)$.

  Relative to an interpretation of prefixes~$P$ we can define:
  \begin{enumerate}
  \item $\mModel{M}$ satisfies $\sFmla{\True}{A}[\sigma]$ iff
    $\mSat{M}{A}[f(\sigma)]$.
  \item $\mModel{M}$ satisfies $\sFmla{\False}{A}[\sigma]$ iff
    $\mSat/{M}{A}[f(\sigma)]$.
  \end{enumerate}
\end{defn}

\begin{defn}
  Let $\Gamma$ be a set of prefixed \olup{formula}s, and let $P(\Gamma)$
  be the set of prefixes that occur in it. If $f$ is an interpretation
  of~$P(\Gamma)$ in $\mModel{M}$, we say that $\mModel{M}$ satisfies
  $\Gamma$ with respect to~$f$, $\mSat{M}{\Gamma}[f]$, if $\mModel{M}$
  satisfies every prefixed \olup{formula} in~$\Gamma$ with respect
  to~$f$. $\Gamma$ is \emph{satisfiable} iff there is a
  model~$\mModel{M}$ and interpretation~$f$ of $P(\Gamma)$ such that
  $\mSat{M}{\Gamma}[f]$.
\end{defn}

\begin{prop}
  If $\Gamma$ contains both $\sFmla{\True}{A}[\sigma]$ and
  $\sFmla{\False}{A}[\sigma]$, for some \olup{formula}~$A$ and
  prefix~$\sigma$, then $\Gamma$ is unsatisfiable.
\end{prop}

\begin{proof}
  There cannot be a model~$\mModel{M}$ and interpretation~$f$ of
  $P(\Gamma)$ such that both $\mSat{M}{A}[f(\sigma)]$ and
  $\mSat/{M}{A}[f(\sigma)]$.
\end{proof}

\begin{thm}[Soundness]
  \ollabel{thm:tableau-soundness}
  If $\Gamma$ has a closed \olup{tableau}, $\Gamma$ is unsatisfiable.
\end{thm}

\begin{proof}
We call a branch of \olup{tableau} satisfiable iff the set of
\olup{signed formula}s on it is satisfiable, and let's call \olup{tableau}
satisfiable if it contains at least one satisfiable branch.

We show the following: Extending a satisfiable \olup{tableau} by one of
the rules of inference always results in a satisfiable \olup{tableau}.
This will prove the theorem: any closed \olup{tableau} results by
applying rules of inference to the \olup{tableau} consisting only of
assumptions from~$\Gamma$. So if $\Gamma$ were satisfiable, any
\olup{tableau} for it would be satisfiable. A closed \olup{tableau},
however, is clearly not satisfiable, since all its branches are closed
and closed branches are unsatisfiable.

Suppose we have a satisfiable \olup{tableau}, i.e., \olup{tableau} with at
least one satisfiable branch. Applying a rule of inference either adds
\olup{signed formula}s to a branch, or splits a branch in two. If the
\olup{tableau} has a satisfiable branch which is not extended by the rule
application in question, it remains a satisfiable branch in the
extended \olup{tableau}, so the extended tableau is satisfiable. So we
only have to consider the case where a rule is applied to a
satisfiable branch.

Let $\Gamma$ be the set of \olup{signed formula}s on that branch, and let
$\sFmla{S}{A}[\sigma] \in \Gamma$ be the \olup{signed formula} to which
the rule is applied. If the rule does not result in a split branch, we
have to show that the extended branch, i.e., $\Gamma$ together with
the conclusions of the rule, is still satisfiable. If the rule results
in split branch, we have to show that at least one of the two
resulting branches is satisfiable.
\tagfalse{prvDiamond}
First, we consider the possible inferences with only one premise.
\begin{enumerate}
\item The branch is expanded by applying $\TRule{\True}{\lnot}$ to
  $\sFmla{\True}{\lnot B}[\sigma] \in \Gamma$. Then the extended
  branch contains the \olup{signed formula}s $\Gamma \cup
  \{\sFmla{\False}{B}[\sigma]\}$.  Suppose $\mSat{M}{\Gamma}[f]$. In
  particular, $\mSat{M}{\lnot B}[f(\sigma)]$. Thus,
  $\mSat/{M}{B}[f(\sigma)]$, i.e., $\mModel{M}$ satisfies
  $\sFmla{\False}{B}[\sigma]$ with respect to~$f$.
\item The branch is expanded by applying $\TRule{\False}{\lnot}$ to
  $\sFmla{\False}{\lnot B}[\sigma] \in \Gamma$: Exercise.
\item The branch is expanded by applying $\TRule{\True}{\land}$ to
  $\sFmla{\True}{B \land C}[\sigma] \in \Gamma$, which results in
  two new \olup{signed formula}s on the branch:
  $\sFmla{\True}{B}[\sigma]$ and $\sFmla{\True}{C}[\sigma]$. Suppose
  $\mSat{M}{\Gamma}[f]$, in particular $\mSat{M}{B \land
    C}[f(\sigma)]$. Then $\mSat{M}{B}[f(\sigma)]$ and
  $\mSat{M}{C}[f(\sigma)]$. This means that $\mModel{M}$ satisfies
  both $\sFmla{\True}{B}[\sigma]$ and $\sFmla{\True}{C}[\sigma]$
  with respect to~$f$.
\item The branch is expanded by applying $\TRule{\False}{\lor}$ to
  $\sFmla{\False}{B \lor C} \in \Gamma$: Exercise.
\item The branch is expanded by applying $\TRule{\False}{\lif}$ to
  $\sFmla{\False}{B \lif C}[\sigma] \in \Gamma$: This results in two
  new \olup{signed formula}s on the branch: $\sFmla{\True}{B}[\sigma]$ and
  $\sFmla{\False}{C}[\sigma]$. Suppose
  $\mSat{M}{\Gamma}[f]$, in particular
  $\mSat/{M}{B \lif C}[f(\sigma)]$. Then
  $\mSat{M}{B}[f(\sigma)]$ and
  $\mSat/{M}{C}[f(\sigma)]$. This means that
  $\mModel{M}, f$ satisfies both
  $\sFmla{\True}{B}[\sigma]$ and $\sFmla{\False}{C}[\sigma]$.

\iftag{prvBox}{%
\item The branch is expanded by applying $\TRule{\True}{\Box}$ to
  $\sFmla{\True}{\Box B}[\sigma] \in \Gamma$:
  \iftag{probBox}{Exercise.}{This results in a new \olup{signed
      formula}~$\sFmla{\True}{B}[\sigma.n]$ on the branch, for some
    $\sigma.n \in P(\Gamma)$ (since $\sigma.n$ must be used).  Suppose
    $\mSat{M}{\Gamma}[f]$, in particular, $\mSat{M}{\Box
      B}[f(\sigma)]$. Since $f$ is an interpretation of prefixes and
    both $\sigma$, $\sigma.n \in P(\Gamma)$, we know that
    $Rf(\sigma)f(\sigma.n)$.  Hence, $\mSat{M}{B}[f(\sigma.n)]$,
    i.e., $\mModel{M}, f$ satisfies $\sFmla{\True}{B}[\sigma.n]$.}

\item The branch is expanded by applying $\TRule{\False}{\Box}$ to
  $\sFmla{\False}{\Box B}[\sigma] \in \Gamma$:
  \iftag{probBox}{Exercise.}{This results in a new \olup{signed
      formula}~$\sFmla{\False}{A}[\sigma.n]$, where $\sigma.n$ is a
    new prefix on the branch, i.e., $\sigma.n \notin P(\Gamma)$.
    Since $\Gamma$ is satisfiable, there is a $\Struct{M}$ and
    interpretation~$f$ of $P(\Gamma)$ such that $\Sat{M}{\Gamma}[f]$,
    in particular $\mSat/{M}{\Box B}[f(\sigma)]$. We have to show
    that $\Gamma \cup \{\sFmla{\False}{B}[\sigma.n]\}$ is
    satisfiable. To do this, we define an interpretation of $P(\Gamma)
    \cup \{\sigma.n\}$ as follows:

  Since $\mSat/{M}{\Box B}[f(\sigma)]$, there is a $w \in W$ such that
  $Rf(\sigma)w$ and $\mSat/{M}{B}[w]$. Let $f'$ be like $f$, except
  that $f'(\sigma.n) = w$. Since $f'(\sigma) = f(\sigma)$ and
  $Rf(\sigma)w$, we have $Rf'(\sigma)f'(\sigma.n)$, so $f'$ is an
  interpretation of $P(\Gamma) \cup \{\sigma.n\}$.  Obviously
  $\mSat/{M}{B}[f'(\sigma.n)]$. Since $f(\sigma') = f'(\sigma')$ for
  all prefixes $\sigma' \in P(\Gamma)$, $\mSat{M}{\Gamma}[f']$.  So,
  $\mModel{M}, f'$ satisfies $\Gamma \cup
  \{\sFmla{\False}{B}[\sigma.n]\}$.}
  }{}

\iftag{prvDiamond}{%
\item The branch is expanded by applying $\TRule{\True}{\Diamond}$ to
  $\sFmla{\True}{\Diamond B}[\sigma] \in \Gamma$:
  \iftag{probDiamond}{Exercise.}{This results in a new \olup{signed
      formula}~$\sFmla{\True}{A}[\sigma.n]$, where $\sigma.n$ is a
    new prefix on the branch, i.e., $\sigma.n \notin P(\Gamma)$.
    Since $\Gamma$ is satisfiable, there is a $\Struct{M}$ and
    interpretation~$f$ of $P(\Gamma)$ such that $\Sat{M}{\Gamma}[f]$,
    in particular $\mSat{M}{\Diamond B}[f(\sigma)]$. We have to show
    that $\Gamma \cup \{\sFmla{\True}{B}[\sigma.n]\}$ is
    satisfiable. To do this, we define an interpretation of $P(\Gamma)
    \cup \{\sigma.n\}$ as follows:

  Since $\mSat{M}{\Diamond B}[f(\sigma)]$, there is a $w \in W$ such that
  $Rf(\sigma)w$ and $\mSat{M}{B}[w]$. Let $f'$ be like $f$, except
  that $f'(\sigma.n) = w$. Since $f'(\sigma) = f(\sigma)$ and
  $Rf(\sigma)w$, we have $Rf'(\sigma)f'(\sigma.n)$, so $f'$ is an
  interpretation of $P(\Gamma) \cup \{\sigma.n\}$.  Obviously
  $\mSat{M}{B}[f'(\sigma.n)]$. Since $f(\sigma') = f'(\sigma')$ for
  all prefixes $\sigma' \in P(\Gamma)$, $\mSat{M}{\Gamma}[f']$.  So,
  $\mModel{M}, f'$ satisfies $\Gamma \cup
  \{\sFmla{\True}{B}[\sigma.n]\}$.}
\item The branch is expanded by applying $\TRule{\False}{\Diamond}$ to
  $\sFmla{\False}{\Diamond B}[\sigma] \in \Gamma$:
  \iftag{probDiamond}{Exercise.}{This results in a new \olup{signed
      formula}~$\sFmla{\False}{B}[\sigma.n]$ on the branch, for some
    $\sigma.n \in P(\Gamma)$ (since $\sigma.n$ must be used).  Suppose
    $\mSat{M}{\Gamma}[f]$, in particular, $\mSat/{M}{\Diamond
      B}[f(\sigma)]$. Since $f$ is an interpretation of prefixes and
    both $\sigma$, $\sigma.n \in P(\Gamma)$, we know that
    $Rf(\sigma)f(\sigma.n)$.  Hence, $\mSat/{M}{B}[f(\sigma.n)]$,
    i.e., $\mModel{M}, f$ satisfies $\sFmla{\False}{B}[\sigma.n]$.}
  }{}
\end{enumerate}
Now let's consider the possible inferences with two premises.
\begin{enumerate}
\item The branch is expanded by applying $\TRule{\False}{\land}$ to
  $\sFmla{\False}{B \land C}[\sigma] \in \Gamma$, which results in
  two branches, a left one continuing through
  $\sFmla{\False}{B}[\sigma]$ and a right one through
  $\sFmla{\False}{C}[\sigma]$. Suppose $\mSat{M}{\Gamma}[f]$, in
  particular $\mSat/{M}{B \land C}[f(\sigma)]$.  Then
  $\mSat/{M}{B}[f(\sigma)]$ or $\mSat/{M}{C}[f(\sigma)]$. In the
  former case, $\mModel{M}, f$ satisfies $\sFmla{\False}{B}[\sigma]$,
  i.e., the left branch is satisfiable. In the latter, $\mModel{M}, f$
  satisfies $\sFmla{\False}{C}[\sigma]$, i.e., the right branch is
  satisfiable.
\item The branch is expanded by applying $\TRule{\True}{\lor}$ to
    $\sFmla{\True}{B \lor C}[\sigma] \in \Gamma$: Exercise.
\item The branch is expanded by applying $\TRule{\True}{\lif}$ to
    $\sFmla{\True}{B \lif C}[\sigma] \in \Gamma$: Exercise.
\end{enumerate}
\end{proof}

\begin{prob}
Complete the proof of \olref[nml][tab][sou]{thm:tableau-soundness}.
\end{prob}

\begin{cor}
\ollabel{cor:entailment-soundness}
If $\Gamma \Proves A$ then $\Gamma \Entails A$.
\end{cor}

\begin{proof}
  If $\Gamma \Proves A$ then for some $B_1$, \dots, $B_n \in
  \Gamma$, $\Delta = \{\sFmla{\False}{A}[1], \sFmla{\True}{B_1}[1],
  \dots, \sFmla{\True}{B_n}[1]\}$ has a closed \olup{tableau}. We want
  to show that $\Gamma \Entails A$. Suppose not, so for some
  $\mModel{M}$ and $w$, $\mSat{M}{B_i}[w]$ for $i=1$, \dots,~$n$, but
  $\mSat/{M}{A}[w]$. Let $f(1) = w$; then $f$ is an interpretation
  of~$P(\Delta)$ into~$\mModel{M}$, and $\mModel{M}$ satisfies~$\Delta$ with
  respect to~$f$. But by \olref{thm:tableau-soundness}, $\Delta$ is
  unsatisfiable since it has a closed \olup{tableau}, a
  contradiction. So we must have $\Gamma \Proves A$ after all.
\end{proof}

\begin{cor}
\ollabel{cor:weak-soundness}
If $\Proves A$ then $A$ is true in all models.
\end{cor}


% ==== OLP file: content/normal-modal-logic/tableaux/tableaux.tex ====
% Part: normal-modal-logic
% Chapter: axioms-systems



\olchapter{nml}{tab}{Modal \usetoken{P}{tableau}}

\begin{editorial}
  Draft chapter on prefixed tableaux for modal logic. Needs more
  examples, completeness proofs, and discussion of how one can find
  countermodels from unsuccessful searches for closed tableaux.
\end{editorial}

\olimport{introduction}
\olimport{rules-for-K}
\olimport{proofs-in-K}
\olimport{soundness}
\olimport{more-rules}
\olimport{more-soundness}
\olimport{simple-S5}
\olimport{completeness}
\olimport{countermodels}

\OLEndChapterHook



\end{document}
